\documentclass[11pt, oneside, article]{memoir}


\settrims{0pt}{0pt} % page and stock same size
\settypeblocksize{*}{35.9pc}{*} % {height}{width}{ratio}
\setlrmargins{*}{*}{1} % {spine}{edge}{ratio}
\setulmarginsandblock{.98in}{.98in}{*} % height of typeblock computed
\setheadfoot{\onelineskip}{2\onelineskip} % {headheight}{footskip}
\setheaderspaces{*}{1.5\onelineskip}{*} % {headdrop}{headsep}{ratio}
\checkandfixthelayout


\usepackage{amsthm}
\usepackage{mathtools}

\usepackage[inline]{enumitem}
\usepackage{ifthen}
\usepackage[utf8]{inputenc} %allows non-ascii in bib file
\usepackage{xcolor}

\usepackage[backend=biber, backref=true, maxbibnames = 10, style = alphabetic]{biblatex}
\usepackage[bookmarks=true, colorlinks=true, linkcolor=blue!50!black,
citecolor=orange!50!black, urlcolor=orange!50!black, pdfencoding=unicode]{hyperref}
\usepackage[capitalize]{cleveref}

\usepackage{tikz}

\usepackage{amssymb}
\usepackage{newpxtext}
\usepackage[varg,bigdelims]{newpxmath}
\usepackage{mathrsfs}
\usepackage{dutchcal}
\usepackage[scaled=0.95]{helvet}
\DeclareMathAlphabet{\mathsf}{T1}{phv}{m}{n}
\usepackage{fontawesome}
\usepackage{ebproof}
\usepackage{stmaryrd}

% cleveref %
  \newcommand{\creflastconjunction}{, and\nobreakspace} % serial comma
  \crefformat{enumi}{\card#2#1#3}
  \crefalias{chapter}{section}


% biblatex %
  \addbibresource{Library20260419.bib}

% hyperref %
  \hypersetup{final}

% enumitem %
  \setlist{nosep}
  \setlistdepth{6}



% tikz %

  \usetikzlibrary{
  	cd,
  	math,
  	decorations.markings,
		decorations.pathreplacing,
  	positioning,
  	arrows.meta,
  	shapes,
		shadows,
		shadings,
  	calc,
  	fit,
  	quotes,
  	intersections,
  	circuits.logic.US,
  }

  \tikzset{
biml/.tip={Glyph[glyph math command=triangleleft, glyph length=.95ex]},
bimr/.tip={Glyph[glyph math command=triangleright, glyph length=.95ex]},
}

\tikzset{
	tick/.style={postaction={
  	decorate,
    decoration={markings, mark=at position 0.5 with
    	{\draw[-] (0,.4ex) -- (0,-.4ex);}}}
  }
}
\tikzset{
	slash/.style={postaction={
  	decorate,
    decoration={markings, mark=at position 0.5 with
    	{\draw[-] (.3ex,.3ex) -- (-.3ex,-.3ex);}}}
  }
}

% oriented WD style (Schultz--Spivak wiring diagrams)
\tikzset{
	oriented WD/.style={
		every to/.style={out=0,in=180,draw},
    label/.style={
    	font=\everymath\expandafter{\the\everymath\scriptstyle},
      inner sep=0pt,
      node distance=2pt and -2pt},
    semithick,
    node distance=1 and 1,
    decoration={markings, mark=at position \stringdecpos with \stringdec},
    ar/.style={postaction={decorate}},
    execute at begin picture={\tikzset{
    	x=\bbx, y=\bby,
      every fit/.style={inner xsep=\bbx, inner ysep=\bby}}}
    },
    string decoration/.store in=\stringdec,
    string decoration={\arrow{stealth};},
    string decoration pos/.store in=\stringdecpos,
    string decoration pos=.7,
    bbx/.store in=\bbx,
    bbx = 1.5cm,
    bby/.store in=\bby,
    bby = 1.5ex,
    bb port sep/.store in=\bbportsep,
    bb port sep=1.5,
    bb port length/.store in=\bbportlen,
    bb port length=4pt,
    bb penetrate/.store in=\bbpenetrate,
    bb penetrate=0,
    bb min width/.store in=\bbminwidth,
    bb min width=1cm,
    bb rounded corners/.store in=\bbcorners,
    bb rounded corners=2pt,
    bb small/.style={
    	bb port sep=1, bb port length=2.5pt, bbx=.4cm, bb min width=.4cm, bby=.7ex},
    bb/.code 2 args={
    	\pgfmathsetlengthmacro{\bbheight}{\bbportsep * (max(#1,#2)+1) * \bby}
      \pgfkeysalso{draw,minimum height=\bbheight,minimum
       width=\bbminwidth,outer sep=0pt,
         rounded corners=\bbcorners,thick,
         prefix after command={\pgfextra{\let\fixname\tikzlastnode}},
         append after command={\pgfextra{\draw
            \ifnum #1=0{} \else foreach \i in {1,...,#1} {
            	($(\fixname.north west)!{\i/(#1+1)}!(\fixname.south west)$) +(-\bbportlen,0) coordinate (\fixname_in\i) -- +(\bbpenetrate,0) coordinate (\fixname_in\i')}\fi
            \ifnum #2=0{} \else foreach \i in {1,...,#2} {
            	($(\fixname.north east)!{\i/(#2+1)}!(\fixname.south east)$) +(-\bbpenetrate,0) coordinate (\fixname_out\i') -- +(\bbportlen,0) coordinate (\fixname_out\i)}\fi;
           }}}
		},
		bb name/.style={
     	append after command={
				\pgfextra{\node[anchor=north] at (\fixname.north) {#1};}
			}
		}
}

\newcommand{\upp}{\begin{tikzcd}[row sep=6pt]~\\~\ar[u, bend left=50pt, looseness=1.3, start anchor=east, end anchor=east]\end{tikzcd}}

\newcommand{\bito}[1][]{
	\begin{tikzcd}[ampersand replacement=\&, cramped]\ar[r, biml-bimr, "#1"]\&~\end{tikzcd}
}
\newcommand{\bifrom}[1][]{
	\begin{tikzcd}[ampersand replacement=\&, cramped]\ar[r, bimr-biml, "{#1}"]\&~\end{tikzcd}
}
\newcommand{\bifromlong}[2][]{
	\begin{tikzcd}[ampersand replacement=\&, column sep=#2, cramped]\ar[r, bimr-biml, "#1"]\&~\end{tikzcd}
}

% Adjunctions
\newcommand{\adj}[5][30pt]{%[size] Cat L, Left, Right, Cat R.
\begin{tikzcd}[ampersand replacement=\&, column sep=#1]
  #2\ar[r, shift left=6pt, "#3"]
  \ar[r, phantom, "\scriptstyle\Rightarrow"]\&
  #5\ar[l, shift left=6pt, "#4"]
\end{tikzcd}
}

\newcommand{\adjr}[5][30pt]{%[size] Cat R, Right, Left, Cat L.
\begin{tikzcd}[ampersand replacement=\&, column sep=#1]
  #2\ar[r, shift left=7pt, "#3"]\&
  #5\ar[l, shift left=7pt, "#4"]
  \ar[l, phantom, "\scriptstyle\Leftarrow"]
\end{tikzcd}
}

\newcommand{\xtickar}[1]{\begin{tikzcd}[baseline=-0.5ex,cramped,sep=small,ampersand
replacement=\&]{}\ar[r,tick, "{#1}"]\&{}\end{tikzcd}}
\newcommand{\xslashar}[1]{\begin{tikzcd}[baseline=-0.5ex,cramped,sep=small,ampersand
replacement=\&]{}\ar[r,slash, "{#1}"]\&{}\end{tikzcd}}



  % amsthm + aliascnt (so cleveref prints the correct environment name) %
\usepackage{aliascnt}

\theoremstyle{definition}
\newtheorem{definitionx}{Definition}[section]

\newaliascnt{examplex}{definitionx}
\newtheorem{examplex}[examplex]{Example}
\aliascntresetthe{examplex}
\crefname{examplex}{Example}{Examples}

\newaliascnt{remarkx}{definitionx}
\newtheorem{remarkx}[remarkx]{Remark}
\aliascntresetthe{remarkx}
\crefname{remarkx}{Remark}{Remarks}

\newaliascnt{notation}{definitionx}
\newtheorem{notation}[notation]{Notation}
\aliascntresetthe{notation}
\crefname{notation}{Notation}{Notations}

\newaliascnt{conventionx}{definitionx}
\newtheorem{conventionx}[conventionx]{Convention}
\aliascntresetthe{conventionx}
\crefname{conventionx}{Convention}{Conventions}

\theoremstyle{plain}

\newaliascnt{theorem}{definitionx}
\newtheorem{theorem}[theorem]{Theorem}
\aliascntresetthe{theorem}
\crefname{theorem}{Theorem}{Theorems}

\newaliascnt{proposition}{definitionx}
\newtheorem{proposition}[proposition]{Proposition}
\aliascntresetthe{proposition}
\crefname{proposition}{Proposition}{Propositions}

\newaliascnt{corollary}{definitionx}
\newtheorem{corollary}[corollary]{Corollary}
\aliascntresetthe{corollary}
\crefname{corollary}{Corollary}{Corollaries}

\newaliascnt{lemma}{definitionx}
\newtheorem{lemma}[lemma]{Lemma}
\aliascntresetthe{lemma}
\crefname{lemma}{Lemma}{Lemmas}

\newaliascnt{warning}{definitionx}
\newtheorem{warning}[warning]{Warning}
\aliascntresetthe{warning}
\crefname{warning}{Warning}{Warnings}

\newaliascnt{fact}{definitionx}
\newtheorem{fact}[fact]{Fact}
\aliascntresetthe{fact}
\crefname{fact}{Fact}{Facts}

\newtheorem*{theorem*}{Theorem}
\newtheorem*{proposition*}{Proposition}
\newtheorem*{corollary*}{Corollary}
\newtheorem*{lemma*}{Lemma}
\newtheorem*{warning*}{Warning}

\newenvironment{example}
  {\pushQED{\qed}\renewcommand{\qedsymbol}{$\lozenge$}\examplex}
  {\popQED\endexamplex}

 \newenvironment{remark}
  {\pushQED{\qed}\renewcommand{\qedsymbol}{$\lozenge$}\remarkx}
  {\popQED\endremarkx}

  \newenvironment{definition}
  {\pushQED{\qed}\renewcommand{\qedsymbol}{$\lozenge$}\definitionx}
  {\popQED\enddefinitionx}

  \newenvironment{convention}
  {\pushQED{\qed}\renewcommand{\qedsymbol}{$\lozenge$}\conventionx}
  {\popQED\endconventionx}


%-------- Single symbols --------%

\DeclareSymbolFont{stmry}{U}{stmry}{m}{n}
\DeclareMathSymbol\fatsemi\mathop{stmry}{"23}

\DeclareFontFamily{U}{mathx}{\hyphenchar\font45}
\DeclareFontShape{U}{mathx}{m}{n}{
      <5> <6> <7> <8> <9> <10>
      <10.95> <12> <14.4> <17.28> <20.74> <24.88>
      mathx10
      }{}
\DeclareSymbolFont{mathx}{U}{mathx}{m}{n}
\DeclareFontSubstitution{U}{mathx}{m}{n}
\DeclareMathAccent{\widecheck}{0}{mathx}{"71}


%-------- Renewed commands --------%

\renewcommand{\ss}{\subseteq}

%-------- Other Macros --------%


\DeclarePairedDelimiter{\present}{\langle}{\rangle}
\DeclarePairedDelimiter{\pairing}{\langle}{\rangle}
\DeclarePairedDelimiter{\copair}{[}{]}
\DeclarePairedDelimiter{\floor}{\lfloor}{\rfloor}
\DeclarePairedDelimiter{\ceil}{\lceil}{\rceil}
\DeclarePairedDelimiter{\corners}{\ulcorner}{\urcorner}
\DeclarePairedDelimiter{\ihom}{[}{]}

\DeclareMathOperator{\Hom}{Hom}
\DeclareMathOperator{\Aut}{Aut}
\DeclareMathOperator{\End}{End}
\DeclareMathOperator{\Mor}{Mor}
\DeclareMathOperator{\dom}{dom}
\newcommand{\Int}{\mathop{\Fun{Int}}\nolimits}
\newcommand{\Lie}{\mathop{\Fun{Lie}}\nolimits}
\DeclareMathOperator{\cod}{cod}
\DeclareMathOperator{\idy}{idy}
\DeclareMathOperator*{\colim}{colim}
\DeclareMathOperator{\im}{im}
\DeclareMathOperator{\ob}{Ob}
\DeclareMathOperator{\Tr}{Tr}
\DeclareMathOperator{\el}{El}




\newcommand{\const}[1]{\texttt{#1}}%a constant, or named element of a set
\newcommand{\Set}[1]{\mathsf{#1}}%a named set
\newcommand{\ord}[1]{\mathsf{#1}}%an ordinal
\newcommand{\cat}[1]{\mathcal{#1}}%a generic category
\newcommand{\Cat}[1]{\mathbf{#1}}%a named category
\newcommand{\fun}[1]{\mathrm{#1}}%a function
\newcommand{\Fun}[1]{\mathsf{#1}}%a named functor




\newcommand{\id}{\mathrm{id}}
\newcommand{\then}{\mathbin{\fatsemi}}

\newcommand{\cocolon}{:\!}


\newcommand{\iso}{\cong}
\newcommand{\too}{\longrightarrow}
\newcommand{\tto}{\rightrightarrows}
\newcommand{\To}[2][]{\xrightarrow[#1]{#2}}
\renewcommand{\Mapsto}[1]{\xmapsto{#1}}
\newcommand{\Tto}[3][13pt]{\begin{tikzcd}[sep=#1, cramped, ampersand replacement=\&, text height=1ex, text depth=.3ex]\ar[r, shift left=2pt, "#2"]\ar[r, shift right=2pt, "#3"']\&{}\end{tikzcd}}
\newcommand{\Too}[1]{\xrightarrow{\;\;#1\;\;}}
\newcommand{\from}{\leftarrow}
\newcommand{\fromm}{\longleftarrow}
\newcommand{\ffrom}{\leftleftarrows}
\newcommand{\From}[1]{\xleftarrow{#1}}
\newcommand{\Fromm}[1]{\xleftarrow{\;\;#1\;\;}}
\newcommand{\surj}{\twoheadrightarrow}
\newcommand{\inj}{\rightarrowtail}
\newcommand{\wavyto}{\rightsquigarrow}
\newcommand{\lollipop}{\multimap}
\newcommand{\imp}{\Rightarrow}
\renewcommand{\iff}{\Leftrightarrow}
\newcommand{\down}{\mathbin{\downarrow}}
\newcommand{\fromto}{\leftrightarrows}
\newcommand{\tickar}{\xtickar{}}
\newcommand{\slashar}{\xslashar{}}
\newcommand{\card}{\,^{\#}}
\newcommand{\bk}[1]{\overset{\scriptscriptstyle\leftarrow}{#1}}


\newcommand{\inv}{^{-1}}
\newcommand{\op}{^\tn{op}}

\newcommand{\tn}[1]{\textnormal{#1}}
\newcommand{\ol}[1]{\overline{#1}}
\newcommand{\ul}[1]{\underline{#1}}
\newcommand{\wt}[1]{\widetilde{#1}}
\newcommand{\wh}[1]{\widehat{#1}}
\newcommand{\wc}[1]{\widecheck{#1}}
\newcommand{\ubar}[1]{\underaccent{\bar}{#1}}



\newcommand{\bb}{\mathbb{B}}
\newcommand{\cc}{\mathbb{C}}
\newcommand{\nn}{\mathbb{N}}
\newcommand{\pp}{\mathbb{P}}
\newcommand{\qq}{\mathbb{Q}}
\newcommand{\zz}{\mathbb{Z}}
\newcommand{\rr}{\mathbb{R}}


\newcommand{\finset}{\Cat{Fin}}
\newcommand{\smset}{\Cat{Set}}
\newcommand{\smcat}{\Cat{Cat}}
\newcommand{\catsharp}{\Cat{Cat}^{\sharp}}
\newcommand{\en}{\Cat{End}}
\newcommand{\vect}{\Cat{Vect}}

\newcommand{\List}{\Fun{List}}
\newcommand{\set}{\tn{-}\Cat{Set}}

\newcommand{\act}{\tn{act}}
\newcommand{\upd}{\tn{upd}}


\newcommand{\yon}{\mathcal{y}}
\newcommand{\poly}{\Cat{Poly}}
\newcommand{\polycart}{\poly^{\tn{cart}}}
\newcommand{\tri}{\mathbin{\triangleleft}}

% lenses
\newcommand{\biglens}[2]{
     \begin{bmatrix}{\vphantom{f_f^f}#2} \\ {\vphantom{f_f^f}#1} \end{bmatrix}
}
\newcommand{\littlelens}[2]{
     \begin{bsmallmatrix}{\vphantom{f}#2} \\ {\vphantom{f}#1} \end{bsmallmatrix}
}
\newcommand{\lens}[2]{
  \relax\if@display
     \biglens{#1}{#2}
  \else
     \littlelens{#1}{#2}
  \fi
}


\newcommand{\slogan}[1]{\begin{center}\textit{#1}\end{center}}

\newcommand{\qand}{\quad\text{and}\quad}
\newcommand{\qqand}{\qquad\text{and}\qquad}
\newcommand{\where}{\qquad\text{where}\quad}

\newcommand{\coto}{\nrightarrow}
\newcommand{\cofun}{{\raisebox{2pt}{\resizebox{2.5pt}{2.5pt}{$\setminus$}}}}

\newcommand{\coalg}{\tn{-}\Cat{Coalg}}

\newcommand{\org}{{\mathbb{O}\Cat{rg}}}

\newcommand{\idcoalg}[1]{{\{\id_{#1}\}}}

\newcommand{\dnote}[1]{{\quad \color{blue}$\lozenge$\;David says:}~#1\;{\color{blue}$\lozenge$}\quad}

%-------- Paper-specific macros --------%

\newcommand{\mfd}{\Cat{Mfd}}
\newcommand{\plenspot}{\Cat{PLensPot}}
\newcommand{\boxx}{\Fun{Box}}
\RenewDocumentCommand{\cot}{g}{\IfValueTF{#1}{\fun{cot}(#1)}{\fun{cot}}}
\newcommand{\ev}{\fun{ev}}
\newcommand{\coev}{\fun{coev}}
\newcommand{\pr}{\fun{pr}}
\newcommand{\Ad}{\fun{Ad}}
\newcommand{\ad}{\fun{ad}}
\newcommand{\lie}[1]{\mathfrak{#1}}
\newcommand{\tpow}[1]{^{\otimes #1}}
\newcommand{\potlens}{\Cat{PotLens}}
\newcommand{\Para}{\Cat{Para}}
\newcommand{\para}[2]{\Para_{#1}(#2)}
\newcommand{\pnla}{\Cat{PNLA}}
\newcommand{\pvect}{\Cat{PVect}}
\newcommand{\Lens}[1]{\Cat{Lens}_{#1}}
\newcommand{\lmfd}{\Lens{\mfd}}
\newcommand{\plm}{\Cat{PLM}}
\newcommand{\inpt}[1]{#1^{-}}
\newcommand{\outp}[1]{#1^{+}}
\newcommand{\lensob}[1]{\binom{\inpt{#1}}{\outp{#1}}}% vertical lens-object notation; to revert, change to (\inpt{#1},\outp{#1})
\newcommand{\lensobs}[2]{\binom{\inpt{#1}\times\inpt{#2}}{\outp{#1}\times\outp{#2}}}% monoidal product of two lens objects

\newcommand{\dec}{\fun{dec}}

% ---- Changeable document parameters ---- %

\linespread{1.1}
%\allowdisplaybreaks
\setsecnumdepth{subsection}
\settocdepth{subsection}
\setlength{\parindent}{15pt}
\setcounter{tocdepth}{1}



%--------------- Document ---------------%
\begin{document}

\title{Nested Potentials}

\author{David I. Spivak}

\date{\vspace{-.2in}}

\maketitle

\begin{abstract}

\end{abstract}

\tableofcontents

%--------------------------------------------------------------------
\chapter{Introduction}\label{ch.intro}
%--------------------------------------------------------------------

% TODO: write introduction
\emph{[Introduction to be written.]}


%--------------------------------------------------------------------
\chapter{Background}\label{ch.background}
%--------------------------------------------------------------------

\section{Notation and preliminaries}\label{sec.prelim}

We fix some typographic conventions. A generic category is written in calligraphic script, $\cat{C},\cat{M},\cat{A}$, and a specifically-named category in boldface, $\Cat{Set},\Cat{Poly},\Cat{Mfd}$. A function or operation is written in roman, $\fun{ev},\fun{pr}$, and a named functor in sans-serif, $\Fun{T},\Fun{Box}$. Lie algebras use fraktur, $\lie{g},\lie{p}$; number systems use blackboard bold, $\nn,\rr$, as does the only bicategory we work with, $\org$. 

In a symmetric monoidal category $(\cat{C},I,\otimes)$, a \emph{$(\otimes)$-comonoid} is an object $c$ equipped with a coassociative comultiplication $\delta_c\colon c\to c\otimes c$ and a counit $\varepsilon_c\colon c\to I$ satisfying the counit law; a \emph{$(\otimes)$-monoid} is the dual. A \emph{comonoid homomorphism} $f\colon c\to c'$ satisfies $f\then\varepsilon_{c'}=\varepsilon_c$ and $f\then\delta_{c'}=\delta_c\then(f\otimes f)$. 

A \emph{supply of comonoids}~\cite{fong2019supplying} is a choice, compatible with $\otimes$, of a cocommutative such structure on every object of $\cat{C}$. Any cartesian monoidal category canonically supplies its diagonals and terminal maps, and every morphism is then automatically a comonoid homomorphism~\cite{fox1976coalgebras}.

\section{Manifolds and paired vector spaces}\label{sec.lie_pnla_bg}

\subsection{Manifolds and notation}\label{sec.manifolds_notation}

We work in the cartesian category $(\mfd,\rr^0,\times)$ of finite-dimensional smooth real manifolds and smooth maps. For a manifold $M$ and point $p\colon\rr^0\to M$, write $T_pM$ for the tangent space at $p$ and $TM$ for the tangent bundle; $T^*_pM$ and $T^*M$ for the cotangent space and bundle. For a smooth map $f\colon M\to N$, write
\[
Tf\colon TM\to TN
\]
for the \emph{tangent map} (the global differential), and
\[
df|_p\colon T_pM\to T_{f(p)}N
\]
for the \emph{differential at $p$}, the restriction of $Tf$ to the fiber over a point $p\in M$. The tangent functor preserves products, so $T(M\times N)\cong TM\times TN$; in particular $T_{(p,q)}(M\times N)\cong T_pM\oplus T_qN$ and dually $T^*_{(p,q)}(M\times N)\cong T^*_pM\oplus T^*_qN$.

We often elide the inclusion $\vect\to\mfd$, treating finite-dimensional real vector spaces as manifolds. The direct sum $(0,\oplus)$ structure in $\vect$ is (a coproduct and) a product, and the inclusion is product preserving. By standard manifold theory, the canonical global chart on a finite-dimensional real vector space $V$ provides linear isomorphisms
\begin{equation}\label{eqn.tangent_vec}
T_pV \;\cong\; V \qqand T^*_pV \;\cong\; V^*
\end{equation}
at every point $p\in V$; we use these identifications throughout, often without comment. In particular, the differential at $p$ of a smooth function $f\colon V\to\rr$ is an element $df|_p\in V^*$, also called \emph{the differential of $f$ at $p$}.
We also use the canonical finite-dimensional dual-sum identification
\begin{equation}\label{eqn.canonical_dual_sum}
(V\oplus V^*)^*\cong V^*\oplus V,\qquad
\lambda\leftrightarrow(\alpha,v)\quad\text{where}\quad
\lambda(v',\alpha')=\alpha(v')+\alpha'(v).
\end{equation}

\subsection{Paired vector spaces}\label{sec.pnla}

\begin{definition}\label{def.pnla}
A \emph{paired vector space} is a finite-dimensional real vector space $V$ together with a linear isomorphism $\sharp_V\colon V^*\To{\cong}V$, the \emph{sharp} map, with inverse $\flat_V\coloneqq\sharp_V\inv\colon V\to V^*$, the \emph{flat} map. The corresponding nondegenerate bilinear form on $V$, the \emph{pairing}, is
\begin{equation}\label{eqn.pairing}
\pairing{v,v'}_V\coloneqq\flat_V(v)(v')\colon V\otimes V\to\rr
\end{equation}
The three forms $\sharp_V$, $\flat_V$, and $\pairing{-,-}_V$ package the same data; we use whichever is most convenient.

We define the category $\pvect$ to have paired vector spaces as objects, and as morphisms $(V,\sharp_V)\to(W,\sharp_W)$ the linear isomorphisms $\phi\colon V\To{\cong}W$ preserving the pairing: $\pairing{\phi v,\phi v'}_W=\pairing{v,v'}_V$ for all $v,v'\in V$.%
\footnote{
Restricting to isomorphisms (making $\pvect$ a groupoid) is essential for our development below: each $\pvect$-morphism is used both as $\phi$ on vectors and as $(\phi^*)\inv$ on covectors, e.g.\ in $T^*\phi$ \eqref{eqn.TT_phi} below.
}
Equivalently, $\phi$ preserves the pairing iff the \emph{pairing-preservation equation} holds:
\begin{equation}\label{eqn.pairing_triangle}
\phi\circ\sharp_V\circ\phi^*=\sharp_W\colon W^*\to W
\qedhere
\end{equation}
\end{definition}

\begin{proposition}\label{prop.pnla_monoidal}
Direct sum gives $\pvect$ a symmetric monoidal structure $(0,\oplus)$, with the zero space (vacuous pairing) as unit and direct-sum pairing $\sharp_{V\oplus W}\coloneqq\sharp_V\oplus\sharp_W$ on $V\oplus W$, equivalently the bilinear form
\begin{equation}\label{eqn.directsum_pairing}
\pairing{(v,w),(v',w')}_{V\oplus W} = \pairing{v,v'}_V+\pairing{w,w'}_W.
\end{equation}
\end{proposition}

\begin{proof}
A direct sum of linear isomorphisms is a linear isomorphism, and \eqref{eqn.directsum_pairing} is direct from $\flat_{V\oplus W}=\flat_V\oplus\flat_W$. Coherence is inherited from $(\vect,0,\oplus)$.\qedhere
\end{proof}

\begin{proposition}\label{prop.pnla_smc_action}
The composite $\pvect\to\vect\hookrightarrow\mfd$ is strong symmetric monoidal, so $\pvect$ acts on $\mfd$ via $((V,\sharp_V),M)\mapsto V\times M$.
\end{proposition}

\begin{proof}
The forgetful $\pvect\to\vect$ preserves direct sums and identities; the inclusion $\vect\hookrightarrow\mfd$ preserves products.\qedhere
\end{proof}

\begin{proposition}\label{prop.canonical_symplectic_pairing}
For any finite-dimensional vector space $V$, the direct sum $V\oplus V^*$ is a paired vector space under the \emph{canonical symplectic pairing}
\begin{equation}\label{eqn.canonical_form}
\pairing{(v,\alpha),(v',\alpha')}_{V\oplus V^*} \;=\; \alpha'(v)-\alpha(v').
\end{equation}
\end{proposition}

\begin{proof}
The form in \eqref{eqn.canonical_form} is bilinear. To see that it is nondegenerate, suppose $(v,\alpha)$ pairs to $0$ with every $(v',\alpha')$. Taking $v'=0$ gives $\alpha'(v)=0$ for all $\alpha'\in V^*$, hence $v=0$; taking $\alpha'=0$ gives $\alpha(v')=0$ for all $v'\in V$, hence $\alpha=0$.\qedhere
\end{proof}

For the canonical symplectic pairing, under \eqref{eqn.canonical_dual_sum}, the associated sharp map is
\begin{equation}\label{eqn.canonical_sharp}
\sharp_{V\oplus V^*}\colon V^*\oplus V\to V\oplus V^*,\qquad(\alpha,v)\mapsto(v,-\alpha).
\end{equation}

\begin{remark}[Sign convention]\label{rmk.swap_pairing}
The bilinear form
\[
\pairing{(v,\alpha),(v',\alpha')}=\alpha'(v)+\alpha(v')
\]
also makes $V\oplus V^*$ a paired vector space; under \eqref{eqn.canonical_dual_sum}, its sharp map is the swap $(\alpha,v)\mapsto(v,\alpha)$ whereas \eqref{eqn.canonical_form} has sharp map \eqref{eqn.canonical_sharp}. We use \eqref{eqn.canonical_form} because its antisymmetry is required for Hamiltonian dynamics (\cref{prop.hamilton,rmk.antisymmetry_substantive}).\qedhere
\end{remark}

\subsection{Tangent and cotangent endofunctors}\label{sec.TT}

For a paired vector space $(V,\sharp_V):\pvect$, define
\begin{equation}\label{eqn.TT_def}
TV\coloneqq V\oplus V \qqand T^*V\coloneqq V\oplus V^*.
\end{equation}
Under the inclusion functor $\vect\to\mfd$ and the global-chart identifications \eqref{eqn.tangent_vec}, these definitions agree with the tangent and cotangent bundles of $V$ as a manifold.

Each is naturally a paired vector space: $TV$ via $\sharp_V\oplus\sharp_V$ (\cref{prop.pnla_monoidal}), $T^*V$ via the canonical symplectic pairing, whose sharp map is \eqref{eqn.canonical_sharp}.
For a $\pvect$-morphism $\phi\colon V\to W$, define
\begin{equation}\label{eqn.TT_phi}
T\phi\coloneqq\phi\oplus\phi
\qqand
 T^*\phi\coloneqq\phi\oplus(\phi^*)\inv.
\end{equation}

\begin{proposition}\label{prop.TT_endofunctors}
The assignments \eqref{eqn.TT_def}, \eqref{eqn.TT_phi} define endofunctors $T,T^*\colon\pvect\to\pvect$.
\end{proposition}

\begin{proof}
If $\phi\colon V\to W$ is a linear isomorphism so are $T\phi$ and $T^*\phi$. Functoriality is obvious. We first check that $T\phi$ preserves pairings.
\begin{align*}
\pairing{T\phi(v_1,v_1'),T\phi(v_2,v_2')}_{TW}
&= \pairing{(\phi v_1,\phi v_1'),(\phi v_2,\phi v_2')}_{TW} && \text{[\eqref{eqn.TT_phi}]}\\
&= \pairing{\phi v_1,\phi v_2}_W+\pairing{\phi v_1',\phi v_2'}_W && \text{[\eqref{eqn.directsum_pairing}]}\\
&= \pairing{v_1,v_2}_V+\pairing{v_1',v_2'}_V && \text{[$\phi$ preserves pairings]}\\
&= \pairing{(v_1,v_1'),(v_2,v_2')}_{TV}.
\end{align*}

Next we check that $T^*\phi$ preserves pairings. Writing $\psi\coloneqq(\phi^*)\inv\colon V^*\to W^*$, the defining property of $\phi^*$ states that $(\phi^*\beta)(v) = \beta(\phi v)$ for all $\beta\in W^*\text{ and }v\in V.$ Setting $\beta\coloneqq\psi\alpha$ and using $\phi^*\circ\psi=\id_{V^*}$, we note that
\begin{equation}\label{eqn.dualeval}
(\psi\alpha)(\phi v) = (\phi^*(\psi\alpha))(v) = \alpha(v),
\end{equation}
and we can conclude that $T^*\phi$ preserves pairings by calculating
\begin{align*}
\pairing{T^*\phi(v_1,\alpha_1),T^*\phi(v_2,\alpha_2)}_{T^*W}
&= (\psi\alpha_2)(\phi v_1)-(\psi\alpha_1)(\phi v_2) && \text{[\eqref{eqn.TT_phi}, \eqref{eqn.canonical_form}]}\\
&= \alpha_2(v_1)-\alpha_1(v_2) && \text{[\eqref{eqn.dualeval}]}\\
&= \pairing{(v_1,\alpha_1),(v_2,\alpha_2)}_{T^*V} && \text{[\eqref{eqn.canonical_form}]}.
\end{align*}
\end{proof}

\begin{proposition}\label{prop.TT_monoidal}
The endofunctors $T,T^*\colon\pvect\to\pvect$ are strong symmetric monoidal with respect to $(0,\oplus)$.
\end{proposition}

\begin{proof}
The unit comparisons are identities since $T0=0\oplus 0=0$ and $T^*0=0\oplus 0^*=0$. The product comparisons are given by interchange:
\begin{align}
\sigma\colon T(V_1\oplus V_2)&\To{\cong}TV_1\oplus TV_2,\quad ((v_1,v_2),(v_1',v_2'))\mapsto((v_1,v_1'),(v_2,v_2')),\label{eqn.T_swap}\\
\sigma^*\colon T^*(V_1\oplus V_2)&\To{\cong}T^*V_1\oplus T^*V_2,\quad ((v_1,v_2),(\alpha_1,\alpha_2))\mapsto((v_1,\alpha_1),(v_2,\alpha_2))\label{eqn.Tstar_swap}
\end{align}
(using $(V_1\oplus V_2)^*\cong V_1^*\oplus V_2^*$ for $\sigma^*$). Both are linear isomorphisms. To check preservation of pairings, compare with a second tangent element $((w_1,w_2),(w_1',w_2'))$ and a second cotangent element $((w_1,w_2),(\beta_1,\beta_2))$. The two sides for $\sigma$ expand to
\[
\pairing{v_1,w_1}_{V_1}+\pairing{v_1',w_1'}_{V_1}+\pairing{v_2,w_2}_{V_2}+\pairing{v_2',w_2'}_{V_2},
\]
and the two sides for $\sigma^*$ expand via \eqref{eqn.canonical_form} to
\[
\beta_1(v_1)+\beta_2(v_2)-\alpha_1(w_1)-\alpha_2(w_2).
\]
\qedhere
\end{proof}

\begin{proposition}\label{prop.legendre}
There is a monoidal natural isomorphism $\sharp\colon T^*\;\xRightarrow{\cong}\;T$ whose components are given by
\[
\id_V\oplus\sharp_V\colon T^*V\to TV,\qquad (v,\alpha)\mapsto(v,\sharp_V(\alpha)).
\]
\end{proposition}

\begin{proof}
Each component is a linear isomorphism by \cref{def.pnla}; naturality is the pairing-preservation axiom of \cref{def.pnla}; monoidality follows from the direct-sum pairing convention $\sharp_{V_1\oplus V_2}=\sharp_{V_1}\oplus\sharp_{V_2}$ (\cref{prop.pnla_monoidal}).\qedhere
\end{proof}

\begin{remark}[Legendre transform]\label{rmk.legendre}
\Cref{prop.legendre} is the (infinitesimal) \emph{Legendre transform} of classical mechanics: $\flat_V$ sends a velocity to its conjugate momentum, and $\sharp_V$ sends a momentum to the corresponding velocity.\qedhere
\end{remark}

The \emph{exponential} of a paired vector space $(V,\sharp_V)$ is the time-1 flow of the constant left-invariant vector field $\sharp_V(\alpha)$ from base point $v$:
\begin{equation}\label{eqn.flow_cot}
\exp_V\colon T^*V\to V,\qquad (v,\alpha)\longmapsto v+\sharp_V(\alpha).
\end{equation}
It is natural in $\pvect$-isomorphisms by linearity and the pairing-preservation equation \eqref{eqn.pairing_triangle}. We use \eqref{eqn.flow_cot} pointwise as the on-directions bijection of $V\yon^V\cong\cot{V}$ in \cref{prop.pnla_polynomial}.

\begin{remark}[Generalization to nilpotent Lie algebras]\label{rmk.pnla_generalization}
\Cref{sec.pnla,sec.TT} generalize from paired vector spaces (abelian $\lie g$) to \emph{paired nilpotent Lie algebras} (PN Lie Algebras): pairs $(\lie g,\pairing{-,-})$ of a finite-dimensional nilpotent real Lie algebra and a pairing on its underlying vector space. Replacing $\pvect$ by $\pnla$ throughout, \cref{prop.pnla_monoidal,prop.canonical_symplectic_pairing,prop.TT_monoidal,prop.legendre} hold verbatim; the rest pick up Lie-group ingredients:
\begin{itemize}
\item \Cref{prop.pnla_smc_action}: the action becomes $((\lie g,\sharp),M)\mapsto\Int\lie g\times M$, where $\Int\lie g$ is the simply-connected Lie group integrating $\lie g$ (Lie's third theorem~\cite[App.~B]{knapp2002lie}; $\Int V\cong V$ for vector spaces).
\item \Cref{prop.TT_endofunctors}: $TV=V\oplus V$ and $T^*V=V\oplus V^*$ in \eqref{eqn.TT_def} are unchanged as vector spaces but now carry brackets mixing $\ad$ and $\ad^*$ (nilpotent of class $\le 2n-1$ when $\lie g$ has class $n$). Morphism formulas \eqref{eqn.TT_phi} are unchanged; bracket-preservation needs $(\phi^*)\inv\circ\ad^*_X=\ad^*_{\phi X}\circ(\phi^*)\inv$, a transpose of $\phi$'s Lie-algebra-homomorphism property.
\item Exponential \eqref{eqn.flow_cot}: extends to $T^*G\to G$, $(g,\alpha)\mapsto g\cdot\exp_G(\sharp_G\alpha)$ for $G\coloneqq\Int\lie g$, via Malcev's diffeomorphism $\exp_G\colon\lie g\To{\cong}\Int\lie g$~\cite[Thm.~1.127]{knapp2002lie}. The vector-space case is abelian: $\Int\lie g=\lie g$, $\exp_G=\id$, multiplication is addition.\qedhere
\end{itemize}
\end{remark}

\subsection{Hamiltonian vector fields}\label{sec.hamiltonian_vector_fields}

\cref{sec.hamiltonian_vector_fields} is not necessary and can be skimmed on a first reading; it will only be used in discussing the wave equation \cref{ch.spring}. We include it because it allows physicists and category theorists to exchange basic intuitions using the basic concepts and results we've developed.

\begin{definition}[Hamiltonian vector field]\label{def.hamiltonian_vector_field}
Let $V$ be a paired vector space and $H\colon T^*V\to\rr$ a smooth function (the \emph{Hamiltonian}). The \emph{Hamiltonian vector field of $H$} is the composite $X_H$ obtained as follows:
\[
\begin{tikzcd}[column sep=1.8em, row sep=2.4em, ampersand replacement=\&]
T^*V
  \ar[d, dashed, "{X_H}"']
  \ar[r, "{(\id,dH)}"]
\&[15pt] T^*V\oplus(T^*V)^*
  \ar[r, equal, "{\eqref{eqn.TT_def},\,\eqref{eqn.canonical_dual_sum}}"]
\&[15pt] (V\oplus V^*)\oplus(V^*\oplus V)
  \ar[d, "{\id\oplus\sharp_{V\oplus V^*}}"]
\\
T(T^*V)
\& {}
\& (V\oplus V^*)\oplus(V\oplus V^*)
  \ar[ll, equal, "{\eqref{eqn.TT_def}^{-1}}"]
\end{tikzcd}
\]
The first arrow sends $(v,\alpha)$ to $((v,\alpha),dH|_{(v,\alpha)})$. The arrow labeled $\id\oplus\sharp_{V\oplus V^*}$ is the component at $T^*V$ of the Legendre transform $\sharp\colon T^*\Rightarrow T$ (\cref{prop.legendre}); explicitly, it leaves the basepoint fixed and applies the canonical symplectic sharp map $\sharp_{V\oplus V^*}=\sharp_{T^*V}$ of \eqref{eqn.canonical_sharp} to the cotangent-fiber component:
\[
((v,\alpha),(\beta,w))\longmapsto((v,\alpha),(w,-\beta)).
\]
The bottom $\eqref{eqn.TT_def}^{-1}$ arrow repackages the explicit direct sum as $T(T^*V)$. Under \eqref{eqn.canonical_dual_sum}, we write
\begin{equation}\label{eqn.hamilton_differential}
dH|_{(v,\alpha)}=(\partial_v H,\partial_\alpha H)\in V^*\oplus V,
\qquad
dH|_{(v,\alpha)}(v',\alpha')=(\partial_v H)(v')+\alpha'(\partial_\alpha H).
\end{equation}
Thus the composite is
\[
X_H(v,\alpha)=((v,\alpha),(\partial_\alpha H,-\partial_v H))\in T(T^*V).\qedhere
\]
\end{definition}

\begin{proposition}[Hamilton's equations]\label{prop.hamilton}
The integral curves $(v(t),\alpha(t))$ of the Hamiltonian vector field of \cref{def.hamiltonian_vector_field} — those whose velocity is the fiber component of $X_H(v,\alpha)$ — satisfy \emph{Hamilton's equations}
\[
\dot v=\partial_\alpha H,\qquad \dot\alpha=-\partial_v H,
\]
and $H$ is conserved along them: $dH/dt=0$.
\end{proposition}

\begin{proof}
By \cref{def.hamiltonian_vector_field}, the fiber component of $X_H(v,\alpha)$ is $(\partial_\alpha H,-\partial_v H)$. Thus the integral-curve condition is
\begin{equation}\label{eqn.hamilton_integral_curve}
(\dot v,\dot\alpha)=(\partial_\alpha H,-\partial_v H),
\end{equation}
which is Hamilton's equations.

For conservation, the chain rule together with \eqref{eqn.hamilton_differential} and \eqref{eqn.hamilton_integral_curve} gives
\[
\begin{aligned}
\frac{dH}{dt}
&= dH|_{(v,\alpha)}(\dot v,\dot\alpha) \\
&= (\partial_v H)(\dot v)+\dot\alpha(\partial_\alpha H) \\
&= (\partial_v H)(\partial_\alpha H)+(-\partial_v H)(\partial_\alpha H)=0.
\end{aligned}
\]
\end{proof}

\begin{remark}[Antisymmetry is substantive]\label{rmk.antisymmetry_substantive}
Had we used the swap pairing of \cref{rmk.swap_pairing} in place of \eqref{eqn.canonical_form}, the vertical map $\id\oplus\sharp_{V\oplus V^*}$ in \cref{def.hamiltonian_vector_field} would use the swap sharp $(\alpha,v)\mapsto(v,\alpha)$ instead of \eqref{eqn.canonical_sharp}. This would give $\dot\alpha=+\partial_v H$ and $dH/dt=2(\partial_v H)(\partial_\alpha H)$, which is generically nonzero: it represents a gradient flow rather than a Hamiltonian flow. Conservation is precisely what the antisymmetry of \eqref{eqn.canonical_form} buys.\qedhere
\end{remark}

\begin{remark}[Coordinate invariance of $X_H$]\label{rmk.XH_natural}
A $\pvect$-isomorphism $\phi\colon V\To{\cong}W$ induces a linear symplectomorphism $T^*\phi\colon T^*V\To{\cong}T^*W$ (\cref{prop.TT_endofunctors}). For any $H\colon T^*W\to\rr$, the chain rule for $H\circ T^*\phi$ together with naturality of the Legendre transform $\sharp$ (\cref{prop.legendre}) at the morphism $T^*\phi$ give
\[
X_H\circ T^*\phi=T(T^*\phi)\circ X_{H\circ T^*\phi},
\]
so $T^*\phi$ takes integral curves of $X_{H\circ T^*\phi}$ on $T^*V$ to those of $X_H$ on $T^*W$. Hamilton's equations transform covariantly under linear canonical transformations.\qedhere
\end{remark}

\begin{remark}[Decoupling of separable Hamiltonians]\label{rmk.XH_separable}
For Hamiltonians $H_i\colon T^*V_i\to\rr$, the separable Hamiltonian on $T^*(V_1\oplus V_2)$ defined in coordinates by
\[
H((v_1,v_2),(\alpha_1,\alpha_2))\coloneqq H_1(v_1,\alpha_1)+H_2(v_2,\alpha_2)
\]
satisfies $\partial_{v_i}H=\partial_{v_i}H_i$ and $\partial_{\alpha_i}H=\partial_{\alpha_i}H_i$, so by \cref{prop.hamilton} the dynamics decouple:
\[
\dot v_i=\partial_{\alpha_i}H_i,\qquad \dot\alpha_i=-\partial_{v_i}H_i\qquad(i=1,2).
\]
At the level of vector fields, the swaps \eqref{eqn.T_swap}, \eqref{eqn.Tstar_swap} of \cref{prop.TT_monoidal} together with the monoidality of $\sharp$ (\cref{prop.legendre}) factor $X_H$ as $X_{H_1}\oplus X_{H_2}$; non-interacting systems evolve independently on their own factors.\qedhere
\end{remark}

\section{Polynomial functors, comonads, and coalgebras}\label{sec.poly_dynamic_bg}

\subsection{The symmetric monoidal closed category $\poly$}\label{sec.poly}

We briefly review the category $\poly$ of polynomial functors; see~\cite{niu2025polynomial} for a comprehensive treatment.

A \emph{polynomial functor} is an endofunctor $p\colon\smset\to\smset$ of the form $\sum_{i: I}\yon^{A_i}$, where $I$ is the set of \emph{positions} and $A_i$ is the set of \emph{directions} at position $i$. Since $p(1)=\sum_{i:I}1=I$, we use the compact notation $p(1)$ for the position set and $p[i]\coloneqq A_i$ for the direction set at $i: p(1)$, so that
\begin{equation}\label{eqn.poly}
p=\sum_{i: p(1)}\yon^{p[i]}=\sum_{i: p(1)}\;\prod_{d: p[i]}\yon.
\end{equation}
We use juxtaposition for scalar multiplication: $pq\coloneqq p\times q$, and we write $1=\yon^0$. A \emph{monomial} is a polynomial of the form $M\yon^N=\sum_{m\in M}\yon^N$; we refer to $\yon^N$ as a \emph{representable} and to $M=M\yon^0$ as a \emph{constant}.

One can derive from the universal property of coproducts and the Yoneda lemma that natural transformations $\varphi\colon p\to q$ are given by:
\begin{equation}\label{eqn.poly_map}
\poly(p,q)=\prod_{i: p(1)}\;\sum_{j: q(1)}\;\prod_{e: q[j]}\; \sum_{d: p[i]}1.
\end{equation}
The outer $\prod_i\sum_j$ yields an ``forward on-positions'' function $\varphi_1\colon p(1)\to q(1)$, which is the component of $\varphi$ at $1:\smset$, and the inner $\prod_e\sum_d$ yields, for each $i:p(1)$, a ``backward on-directions'' function from the directions of $q$ at $\varphi_1(i)$ to the directions of $p$ at $i$; we denote this function
\[\bk{\varphi}_i\colon q[\varphi_1(i)]\to p[i],\]
writing the overline-bracket $\bk{(\,\cdot\,)}$ for the backward component throughout. We take the morphisms $p\to q$ in $\poly$ to be the natural transformations between the underlying functors.

For polynomials $p,q$, their \emph{Dirichlet product} is
\begin{equation}\label{eqn.dirichlet}
p\otimes q\coloneqq\sum_{(i,j):p(1)\times q(1)}\yon^{p[i]\times q[j]}\;.
\end{equation}
The unit is $\yon=\yon^1$. The Dirichlet product is symmetric monoidal and closed: the \emph{internal hom} is
\begin{equation}\label{eqn.dirichlet_hom}
\ihom{p,q}=\prod_{i: p(1)}\;\sum_{j: q(1)}\;\prod_{e: q[j]}\;\sum_{d: p[i]}\yon
\end{equation}
satisfying $\poly(r\otimes p,\, q)\cong\poly(r,\,\ihom{p,q})$ naturally in $r$. Note that $\ihom{p,q}(1)=\poly(p,q)$, as is immediate from comparing \eqref{eqn.dirichlet_hom} with \eqref{eqn.poly_map}. We can rewrite it as
\[
\ihom{p,q}\cong\sum_{\varphi: \poly(p,q)}\yon^{\sum_{i: p(1)}q[\varphi_1(i)]}
\]
A direction at $\varphi:\ihom{p,q}(1)$ consists of a pair $(i,e)$, where $i:p(1)$ is a $p$-position and $e:q[\varphi_1(i)]$ is a $q$-direction.

\subsection{Polynomial comonads as categories}\label{sec.comonads}

The \emph{substitution} (or \emph{composition}) \emph{product} on $\poly$ is defined by $(p\tri q)(S)\coloneqq p(q(S))$, with unit $\yon$. Using the expanded form of \eqref{eqn.poly}:
\begin{equation}\label{eqn.substitution}
p\tri q=\sum_{i: p(1)}\;\prod_{d: p[i]}\;\sum_{j: q(1)}\;\prod_{e: q[j]}\yon.
\end{equation}
A position of $p\tri q$ is a pair $(i,j)$ where $i: p(1)$ and $j\colon p[i]\to q(1)$; the direction set is $(p\tri q)[(i,j)]=\sum_{d:p[i]}q[j_d]$. In particular, viewing a set $S:\smset$ as the constant polynomial $S\yon^0$, substitution recovers evaluation: $p\tri S=p(S)$.

A \emph{$\tri$-comonoid} in $\poly$ is a comonoid for the substitution product: a polynomial $c$ equipped with a counit $\epsilon\colon c\to\yon$ and comultiplication $\delta\colon c\to c\tri c$ satisfying coassociativity and counitality. Since the substitution product is composition of endofunctors on $\smset$, a $\tri$-comonoid is equivalently a comonad on $\smset$ whose underlying endofunctor is polynomial; we call it a \emph{polynomial comonad}.

Polynomial comonads are---up to isomorphism---small categories. Writing $c=\sum_{i:c(1)}\yon^{c[i]}$, the positions $c(1)$ are the objects of some category $\cat{C}$, and for each object $i:c(1)$, the directions $c[i]$ are the morphisms of $\cat{C}$ with domain $i$:
\[
c[i]\coloneqq\sum_{j:c(1)}\cat{C}(i,j).
\]
The counit $\epsilon\colon c\to\yon$ picks out an identity morphism $\id_i:c[i]$ at each object $i$. The comultiplication $\delta\colon c\to c\tri c$ assigns codomains and composition: given an object $i$ and a morphism $f:c[i]$, the on-positions component of $\delta$ returns $\cod(f):c(1)$, and given a further morphism $g:c[\cod(f)]$, the on-directions component returns the composite $f\then g:c[i]$. The comonoid laws guarantee that $\cod(\id_i)=i$, that composition is unital, that $\cod(f\then g)=\cod(g)$, and that composition is associative.

\begin{proposition}[Polynomial comonads are categories {\cite{ahman2016directed}; see also \cite{niu2025polynomial}}]\label{prop.comonads_cats}
The above correspondence defines an isomorphism between polynomial comonads and small categories.
\end{proposition}

\begin{example}[The store comonad and the pair groupoid]\label{ex.store}
For any set $S$, the monomial $S\yon^S$ arises as the store comonad from the adjunction $S\times-\;\dashv\;(-)^S$ on $\smset$, i.e.\ the comonad with underlying endofunctor $X\mapsto S\times X^S$. We can also write
\begin{equation}\label{eqn.store_decomp}
S\yon^S\;\cong\;\ihom{S\yon,\yon}\otimes S\yon
\end{equation}
since $\ihom{S\yon,\yon}=\yon^S$, and $\yon^S\otimes S\yon\cong S\yon^{S\times 1}\cong S\yon^S$. The counit $\epsilon\colon S\yon^S\to\yon$ sends each position $s:S$ to the identity $\id_s\coloneqq s:S=c[s]$. The comultiplication $\delta\colon S\yon^S\to S\yon^S\tri S\yon^S$ acts as follows: on positions, $\delta$ sends $s:S$ to the pair $(s,\cod)$ where $\cod\colon S\to S$ is the identity---so $\cod(t)=t$ for all $t:S$. On directions, given $s:S$ and $t:c[s]=S$, the codomain is $\cod(t)=t$, and given a further $u:c[t]=S$, the composite is $t\then u\coloneqq u$.

By \cref{prop.comonads_cats} this is a category: it has object set $S$ and, for each pair $(s,t):S\times S$, exactly one morphism $s\to t$. The composition $t\then u=u$ says that composites are determined by the codomain of the second factor, as expected when every hom-set is a singleton. Every arrow $s\to t$ is an isomorphism with inverse $t\to s$, so this is the \emph{pair groupoid} (or \emph{codiscrete category}) on $S$.
\end{example}



\subsection{Polynomial coalgebras as dynamical systems}\label{sec.coalgebras}

For a polynomial $p$, a \emph{$p$-coalgebra} is a pair $(S,\beta)$ where $S:\smset$ is a set and $\beta\colon S\to p(S)$ is a function. We call elements of $S$ states; given a state $s:S$, the map assigns a \emph{readout} position $r_s\coloneqq\beta_1(s):p(1)$ and an \emph{update} function $\bk{\beta}_s\colon p[r_s]\to S$ sending each direction at $r_s$ to a new state. For example when $p=A\yon^B$ then a $p$-coalgebra $S\to AS^B$ is equivalently a readout function $S\to A$ and an update function $S\times B\to S$.

A $p$-coalgebra is thus a deterministic dynamical system whose interface is $p$: the state exposes an output in $p(1)$ and transitions in response to inputs in $p[\beta_1(s)]$, capturing Moore machines, discretized differential equations, Markov decision processes, and other mode-dependent dynamics in a single framework~\cite{niu2025polynomial}. It is a special case of the general theory of coalgebras for endofunctors; see~\cite{jacobs2017introduction} for a comprehensive treatment.

\begin{proposition}[Coalgebras as polynomial maps]\label{prop.coalg_as_poly_map}
For $p:\poly$ and $S:\smset$, there is a bijection
\[
\smset(S,\,p\tri S)\;\cong\;\poly(S\yon^S,\,p),
\]
natural in $p:\poly$ and in $S:\smset_\cong$.
\end{proposition}

\begin{proof}
A coalgebra $\beta\colon S\to p\tri S$ unpacks as a pair $\beta=(\beta_1,\beta^\sharp)\colon S\to\sum_{i:p(1)}\prod_{d:p[i]}S$: for each $s:S$, a position $\beta_1(s):p(1)$ and an update $\beta^\sharp_s\colon p[\beta_1(s)]\to S$. A polynomial map $\hat\beta\colon S\yon^S\to p$ unpacks as a pair $(\hat\beta_1,\hat\beta^\sharp)$ of identical signature by \eqref{eqn.poly_map}. The bijection sets $\hat\beta_1\coloneqq\beta_1$ and $\hat\beta^\sharp\coloneqq\beta^\sharp$.\footnote{
Equivalently, $\poly$ has a left Kan operation $\lens{p}{q}$ with $\poly(p,q'\tri q)\cong\poly(\lens{p}{q},q')$ given by $\lens{p}{q}=\sum_{i:p(1)}\yon^{q\tri p[i]}$, so $\lens{S}{S}=S\yon^S$.
}
Naturality in $p:\poly$ is just composition. For naturality in $S$, a bijection $f\colon S\To{\cong}T$ acts as $f\yon^{f\inv}\colon S\yon^S\to T\yon^T$ (positions $f$, directions $f\inv$), and we need to check
\[
\begin{tikzcd}
	S\ar[r,"\beta"]\ar[d, "f"']&p\tri S\ar[d, "p\tri f"]\\
	T\ar[r,"\beta'"']&p\tri T
\end{tikzcd}
\quad\leftrightsquigarrow\quad
\begin{tikzcd}
	S\yon^S\ar[r,"\hat\beta"]\ar[d, "f\yon^{f\inv}"']&
	p\ar[d, equal]\\
	T\yon^T\ar[r,"\hat{\beta'}"']&
	p
\end{tikzcd}
\]
The first commutes iff $\beta'_1\circ f=\beta_1$ and $f\circ\beta^\sharp_s=\beta'^\sharp_{f(s)}$ for all $s:S$; the second iff $\hat{\beta'}_1\circ f=\hat\beta_1$ and $\hat\beta^\sharp_s=f\inv\circ\hat{\beta'}^\sharp_{f(s)}$. Under the bijection $\hat\beta_\bullet=\beta_\bullet$, these coincide ($f$ being a bijection).\qedhere
\end{proof}

The Dirichlet product formula \eqref{eqn.dirichlet} gives $1\yon^1=\yon$ and $S\yon^S\otimes T\yon^T=(S\times T)\yon^{S\times T}$, so the assignment $S\mapsto S\yon^S$ extends to a strong symmetric monoidal functor $\smset_\cong\to\poly$, sending a bijection $f\colon S\To{\cong}T$ to $f\yon^{f\inv}$ (\cref{prop.coalg_as_poly_map}). Composing with the Dirichlet product gives a strong monoidal action of $\smset_\cong$ on $\poly$.

\begin{definition}[$\Fun{smpl}$ and the store action]\label{def.store_action}
Write
\[
\Fun{smpl}\colon(\smset_\cong,1,\times)\to(\poly,\yon,\otimes),\qquad S\longmapsto S\yon^S
\]
for the above functor, and define the \emph{store action} by
\[
(\cdot)\colon\smset_\cong\times\poly\to\poly
\where
S\cdot p\coloneqq S\yon^S\otimes p.
\qedhere
\]
\end{definition}


\subsection{Cotangent bundles as polynomials}\label{sec.cot}

The category of smooth manifolds provides the source of polynomials relevant to physics. Each manifold has a cotangent bundle, and each smooth map has a Jacobian transpose; combining these gives a functor into $\poly$.

\begin{definition}[Cotangent functor]\label{def.cot}
The \emph{cotangent functor} $\cot\colon\mfd\to\poly$ sends a smooth manifold $M$ to
\begin{equation}\label{eqn.cot_def}
\cot{M}\coloneqq\sum_{m\in M}\yon^{T^*_mM}
\end{equation}
where $T^*_mM$ is the cotangent space at $m$. A smooth map $f\colon M\to N$ is sent to the $\poly$ map $\cot{f}\colon\cot{M}\to\cot{N}$ with forward part $\cot{f}_1=f$ and backward part $\bk{\cot{f}}_m\coloneqq(Df_m)^\top\colon T^*_{f(m)}N\to T^*_mM$ (the Jacobian transpose at $m$).
\end{definition}

This is indeed functorial: the identity map $\id_M\colon M\to M$ has Jacobian $D(\id_M)_m=\id_{T_mM}$, whose transpose is $\id_{T^*_mM}$, so $\cot{\id_M}=\id_{\cot{M}}$. For composable smooth maps $f\colon M\to N$ and $g\colon N\to P$, the chain rule gives $D(g\circ f)_m=Dg_{f(m)}\circ Df_m$, so
\[
(D(g\circ f)_m)^\top=(Df_m)^\top\circ(Dg_{f(m)})^\top=\bk{\cot{f}}_m\circ\bk{\cot{g}}_{f(m)},
\]
which is the backward part of $\cot{f}\then\cot{g}$.

\begin{proposition}\label{prop.cot_monoidal}
The cotangent functor $\cot\colon(\mfd,\mathbb{R}^0,\times)\to(\poly,\yon,\otimes)$ is strong symmetric monoidal:
\[
\cot{\mathbb{R}^0}\cong\yon
\qqand
\cot{M\times N}\cong\cot{M}\otimes\cot{N}.
\]
\end{proposition}

\begin{proof}
\textit{Unit.} The one-point manifold $\mathbb{R}^0=\{*\}$ has $T^*_*\mathbb{R}^0=0$, so $\cot{\mathbb{R}^0}=\sum_{*:1}\yon^0=\yon$.

\textit{Products.} Since $T$ preserves products, the cotangent space of a product splits as $T^*_{(m,n)}(M\times N)\cong T^*_mM\oplus T^*_nN$, hence as an underlying set is canonically $T^*_mM\times T^*_nN$. Thus
\begin{equation}\label{eqn.cot_times}
\cot{M\times N}=\sum_{(m,n)\in M\times N}\yon^{T^*_mM\times T^*_nN}=\cot{M}\otimes\cot{N}
\end{equation}
using the definition of the Dirichlet product \eqref{eqn.dirichlet}.

\textit{Naturality.} For a product map $f\times g\colon M\times N\to M'\times N'$, the Jacobian is block-diagonal: $D(f\times g)_{(m,n)}=Df_m\oplus Dg_n$. Hence $(D(f\times g)_{(m,n)})^\top=(Df_m)^\top\times(Dg_n)^\top$, so the isomorphism $\cot{M\times N}\cong\cot{M}\otimes\cot{N}$ is natural in $M$ and~$N$.
%
%\textit{Coherence.} The associator and unitor for the Dirichlet product correspond to the canonical isomorphisms $(M\times N)\times P\cong M\times(N\times P)$ and $M\times 1\cong M$ in $\mfd$, so the coherence diagrams for a strong monoidal functor commute.
\end{proof}

\Cref{prop.cot_monoidal} uses only the underlying sets of the cotangent fibers, e.g.\ in \eqref{eqn.cot_times}. The linear structure is not lost, however: \cref{lem.cot_comonoid} shows that the diagonal and terminal maps in $\mfd$ induce fiberwise addition and zero on $\cot{M}$.

\begin{remark}
In fact, $\cot$ also preserves coproducts,
\[
\cot{M}+\cot{N}\cong\cot{M+N}.
\]
Hence $\cot\colon(\mfd,0,+,1,\times)\to(\poly,0,+,\yon,\otimes)$ is a distributive monoidal functor.
\end{remark}

\begin{lemma}\label{lem.cot_comonoid}
For any smooth manifold $M$, the polynomial $\cot{M}$ is a cocommutative comonoid in $(\poly,\yon,\otimes)$. The comultiplication and counit
\[
\cot{\Delta_M}\colon\cot{M}\to\cot{M}\otimes\cot{M}
\qqand
 \cot{!_M}\colon\cot{M}\to\yon
\]
sends each position $m\in M$ to $(m,m)$. On directions, each cotangent fibre $T^*_mM$ carries the induced cocommutative comonoid structure
\[
T^*_mM\times T^*_mM\to T^*_mM,\quad (\alpha,\alpha')\mapsto\alpha+\alpha'
\qqand 1\to T^*_mM,\quad *\mapsto 0,
\]
i.e.\ fiberwise addition and zero. Thus the $\otimes$-comonoid structure on $\cot{M}$ recovers the additive structure of each cotangent space.%
\footnote{
One can recover scalar structure on $T^*_mM$ as well. For any $M:\mfd$, there is a natural map
\[
(\cdot)\colon\cot{M}\to\yon^\rr\otimes\cot{M}
\]
Axioms relating the comonoid structure of $\cot{M}$ and various structures on $\rr$ enforce that the forward map is $x\mapsto (!,x)$ and that the backwards map $\bk{(\cdot)}_x\colon (r,\alpha)\mapsto r\cdot\alpha$ encodes the following laws:
\[
	1\cdot\alpha=\alpha,\qquad
	(r_1r_2)\cdot\alpha=r_1\cdot(r_2\cdot\alpha),\qquad
	0\cdot\alpha=0,\qand
	r\cdot(\alpha_1+\alpha_2)=(r\cdot\alpha_1)+(r\cdot\alpha_2).
\]
Notably missing is the axiom $(r_1+r_2)\cdot\alpha=(r_1\cdot\alpha)+(r_2\cdot\alpha)$, which seemingly cannot be encoded this way because there is seemingly no way to copy a cotangent vector ($\alpha$).
}
\end{lemma}

\begin{proof}
Since $(\mfd,\mathbb{R}^0,\times)$ is cartesian, every manifold $M$ is a cocommutative comonoid via the diagonal $\Delta_M\colon M\to M\times M$ and $!_M\colon M\to\mathbb{R}^0$. Since $\cot$ is strong symmetric monoidal (\cref{prop.cot_monoidal}), it preserves comonoid structures. On positions this gives $m\mapsto(m,m)$ and $m\mapsto *$. For the backward maps: the Jacobian of $\Delta_M$ at $m$ is the diagonal $T_mM\to T_mM\times T_mM$, $v\mapsto(v,v)$; its transpose is addition $T^*_mM\times T^*_mM\to T^*_mM$. The Jacobian of $!_M$ at $m$ is $T_mM\to 0$, whose transpose is the map $1\to T^*_mM$ that selects the zero covector.
\end{proof}


\begin{proposition}[Hopf monoid structure on cotangent bundles of Lie groups]\label{prop.cot_hopf}
Let $G$ be a Lie group with unit $\eta\colon\mathbb{R}^0\to G$, multiplication $\mu\colon G\times G\to G$, and inversion $\iota\colon G\to G$. Then $\cot{G}$ is a Hopf monoid in $(\poly,\yon,\otimes)$: the monoid structure is
\[
\cot{\eta}\colon\yon\to\cot{G},\qquad \cot{\mu}\colon\cot{G}\otimes\cot{G}\to\cot{G},
\]
the comonoid structure is as in \cref{lem.cot_comonoid}, and the antipode is $\cot{\iota}\colon\cot{G}\to\cot{G}$.
\end{proposition}

\begin{proof}
The Lie group axioms are equations in $\mfd$; applying the strong monoidal functor $\cot$ preserves them. In particular, the monoid axioms for $(\eta,\mu)$ become monoid axioms for $(\cot{\eta},\cot{\mu})$, the compatibility between multiplication and comultiplication (which holds because $\Delta$ is natural), and the antipode condition $\mu\circ(\iota\times\id)\circ\Delta=\eta\circ!=\mu\circ(\id\times\iota)\circ\Delta$ maps to the Hopf monoid axiom in $\poly$. See also~\cite{spivak2023hopf}.
\end{proof}

Recall the inclusion $\pvect\to\vect\hookrightarrow\mfd$ (\cref{prop.pnla_smc_action}) and $\Fun{smpl}(S)\coloneqq S\yon^S$ from \cref{def.store_action}.

\begin{proposition}\label{prop.pnla_polynomial}
There is a weakly commuting square
\begin{equation}\label{eqn.pnla_polynomial_square}
\begin{tikzcd}[ampersand replacement=\&]
\pvect\ar[r,"\Fun{inc}"]\ar[d,"|{-}|"']
\ar[dr,phantom,"\cong^\theta"description] \& \mfd\ar[d,"\cot"]\\
\smset_\cong\ar[r,"\Fun{smpl}"'] \& \poly
\end{tikzcd}
\end{equation}
The 2-cell $\theta_V\colon V\yon^V\To{\cong}\cot{V}$ at each $V:\pvect$ is identity on positions and, on directions at $v\in V$, the cotangent flow \eqref{eqn.flow_cot}
\begin{equation}\label{eqn.directions_sharp}
T^*_vV\To{\eqref{eqn.tangent_vec}}V^*\To{\sharp_V}V\To{v+(-)}V,
\qquad
\alpha\longmapsto v+\sharp_V(\alpha).
\end{equation}
\end{proposition}

\begin{proof}
On positions, both functors send $V$ to its underlying set, so the on-positions component is the identity. At each $v\in V$, \eqref{eqn.directions_sharp} is the composite of two linear isomorphisms with translation by $v$, hence a bijection. So $\Fun{smpl}(|V|)=V\yon^V\cong\cot{V}$ in $\poly$ for each $V$. For naturality at a $\pvect$-isomorphism $\phi\colon V\to W$, the on-positions square commutes by definition. On directions at $v$, we must check that the two composites $T^*_{\phi(v)}W\to V$ agree:
\begin{align*}
\beta\;&\Mapsto{\phi^*}\;\phi^*\beta\;\Mapsto{\sharp_V}\;\sharp_V(\phi^*\beta)\;\Mapsto{v+(-)}\;v+\sharp_V(\phi^*\beta), && \text{[$F^V_v\circ\cot{\phi}$]}\\
\beta\;&\Mapsto{\sharp_W}\;\sharp_W(\beta)\;\Mapsto{\phi(v)+(-)}\;\phi(v)+\sharp_W(\beta)\;\Mapsto{\phi^{-1}}\;v+\phi^{-1}(\sharp_W\beta) && \text{[$\phi\yon^\phi\circ F^W_{\phi(v)}$]}.
\end{align*}
The two are equal iff $\sharp_V\circ\phi^*=\phi^{-1}\circ\sharp_W$, which is \eqref{eqn.pairing_triangle}.\qedhere
\end{proof}

\begin{remark}\label{rmk.pnla_fiberwise_addition}
By \cref{prop.comonads_cats}, a $\tri$-comonoid structure on $V\yon^V$ is a category whose object set is $V$ and whose arrows out of each object are indexed by $V$. By \cref{ex.store}, the canonical such structure is the codiscrete category on $V$: there is exactly one morphism $v\to v'$ for each pair, and the direction-set $V$ at position $v$ is interpreted as the set of codomains.

Transporting this $\tri$-comonoid structure to $\cot{V}$ along the iso of \cref{prop.pnla_polynomial}, a covector $\alpha\in T^*_vV$ at $v$ corresponds to the unique morphism $v\to v+\sharp_V(\alpha)$ in the codiscrete category on $V$. The identity at $v$ corresponds to the zero covector $0\in T^*_vV$, since $v+\sharp_V(0)=v$. Composition of $\alpha$ at $v$ with $\alpha'$ at $v+\sharp_V(\alpha)$ corresponds to $\alpha+\alpha'$ at $v$, since both produce the morphism with codomain $v+\sharp_V(\alpha)+\sharp_V(\alpha')=v+\sharp_V(\alpha+\alpha')$. Thus the pairing $\sharp_V$ identifies $\cot{V}$ with the codiscrete category on $V$, with codomain map $(v,\alpha)\mapsto v+\sharp_V(\alpha)$ and composition by fiberwise addition of covectors.\qedhere
\end{remark}

\section{Dynamic category-like structures}\label{sec.dynamic_structures}

\subsection{Wiring-diagram operads}\label{sec.wd_operads}

Categories, monoidal categories, traced categories, operads, etc.\ can be understood in terms of the sorts of compositions they allow. For example:
\begin{equation}\label{eqn.wd_examples}
\begin{tikzpicture}[baseline=(Catg)]
\begin{scope}[font=\footnotesize, text height=1.5ex, text depth=.5ex]
  \begin{scope}[oriented WD, bb port sep=1, bb port length=0, bb min width=.4cm, bby=.2cm, inner xsep=.2cm, x=.5cm, y=.3cm]
  	\node[bb={1}{1}] (Catf) {$f_1$};
  	\node[bb={1}{1}, right=1 of Catf] (Catg) {$f_2$};
  	\node[bb={0}{0}, fit=(Catf) (Catg)] (Cat) {};
  	\node[coordinate] at (Cat.west|-Catf_in1) (Cat_in1) {};
  	\node[coordinate] at (Cat.east|-Catg_out1) (Cat_out1) {};
  	\draw (Cat_in1) -- (Catf_in1);
  	\draw (Catf_out1) -- (Catg_in1);
  	\draw (Catg_out1) -- (Cat_out1);
  	\draw[label]
  		node[left=3pt of Cat_in1] {$A$}
  		node at ($(Catf_out1)!.5!(Catg_in1)+(0,.8)$) {$B$}
  		node[right=3pt of Cat_out1] {$C$}
  	;
  %
  	\node[bb={1}{2}, above right=-1 and 6.5 of Catf] (Monf) {$f_1$};
  	\node[bb={2}{1}, below right=-.5 and 1 of Monf] (Mong) {$f_2$};
  	\node[bb={0}{0}, fit=(Monf) (Mong)] (Mon) {};
  	\node[coordinate] at (Mon.west|-Monf_in1) (Mon_in1) {};
  	\node[coordinate] at (Mon.west|-Mong_in2) (Mon_in2) {};
  	\node[coordinate] at (Mon.east|-Monf_out1) (Mon_out1) {};
  	\node[coordinate] at (Mon.east|-Mong_out1) (Mon_out2) {};
  	\draw  (Mon_in1) -- (Monf_in1);
  	\draw (Monf_out1) -- (Mon_out1);
  	\draw (Monf_out2) to (Mong_in1);
  	\draw (Mon_in2) -- (Mong_in2);
  	\draw (Mong_out1) -- (Mon_out2);
  	\draw[label]
  		node[left=3pt of Mon_in1] {$A$}
  		node[left=3pt of Mon_in2] {$B$}
  		node[right=3pt of Mon_out1] {$C$}
  		node at ($(Monf_out2)!.5!(Mong_in1)+(.2,.4)$) {$D$}
  		node[right=3pt of Mon_out2] {$E$}
  	;
  %
  	\node[bb={2}{1}, below right= -1.75 and 6.5 of Monf] (Trf) {$f_1$};
  	\node[bb={2}{2}, below right=-1 and 1 of Trf] (Trg) {$f_2$};
  	\node[bb={0}{0}, fit={($(Trf.north west)+(-.5,2)$) ($(Trg.south east)+(.5,0)$)}] (Tr) {};
  	\node[coordinate] at (Tr.west|-Trf_in2) (Tr_in1) {};
  	\node[coordinate] at (Tr.west|-Trg_in2) (Tr_in2) {};
  	\node[coordinate] at (Tr.east|-Trf_out1) (Tr_out1) {};
  	\node[coordinate] at (Tr.east|-Trg_out2) (Tr_out2) {};
  	\draw (Tr_in1) -- (Trf_in2);
  	\draw (Trf_out1) to (Trg_in1);
  	\draw (Tr_in2) -- (Trg_in2);
  	\draw (Trg_out2) -- (Tr_out2);
  	\draw let \p1=(Trg.east), \p2=(Trf.north west), \n1=\bbportlen, \n2=\bby in
  		(Trg_out1) to[in=0] (\x1+\n1,\y2+\n2) -- (\x2-\n1,\y2+\n2) to[out=180] (Trf_in1);
  	\draw[label]
  		node[left=3pt of Tr_in1] {$A$}
  		node[left=3pt of Tr_in2] {$B$}
  		node at ($(Trf_out1)!.4!(Trg_in1)+(0,1)$) {$D$}
  		node[right=3pt of Tr_out2] {$E$}
  		node at ($(Trf.north)+(0,1.5)$) {$C$}
  	;
  \end{scope}
  \begin{scope}[circuit logic US, thick, every to/.style={out=0,in=180}]
  	\node[and gate, draw, above right=-.2 and 3 of Trf] (Opdf) {$f_1$};
  	\node[and gate, draw, below right=-.1 and 0.5 of Opdf] (Opdg) {$f_2$};
  	\node[and gate, inner ysep=-4pt, draw, fit={($(Opdf.north west)+(0,-.2)$) ($(Opdg.south east)+(0,.2)$)}] (Opd) {};
  	\node[coordinate] at (Opd.input 1|-Opdf.west) (Opd_in1) {};
  	\node[coordinate] at (Opd.input 1|-Opdg.input 2) (Opd_in2) {};
  	\node[coordinate] at (Opd.output) (Opd_out) {};
  	\draw (Opd_in1) to (Opdf.west);
  	\draw (Opd_in2) to (Opdg.input 2);
  	\draw (Opdf.output) to (Opdg.input 1);
  	\draw (Opdg.output) to (Opd_out);
  	\draw[label]
  		node[left=3pt of Opd_in1] {$A$}
  		node[left=3pt of Opd_in2] {$B$}
  		node at ($(Opdf.output)!.5!(Opdg.input 1)+(.1,.25)$) {$C$}
  		node[right=3pt of Opd_out] {$D$}
  	;
  \end{scope}
	\node[below=.3 of Cat.south] (Cat name) {category};
	\node[text width=1.5cm, align=center] at (Cat name-|Mon) {monoidal\\category};
	\node[text width=2.5cm, align=center] at (Cat name-|Tr) {traced monoidal\\category};
	\node at (Cat name-|Opd) {operad};
\end{scope}
\end{tikzpicture}
\end{equation}
\smallskip

\noindent In each case there is an associated operad $\cat{W}$ whose objects are the legal interfaces (boxes-with-objet-labeled-ports) and whose morphisms $\phi\colon b_1,\ldots,b_K\to c$ are the legal arrangements of boxes in a box. That is, there are operads $\cat{W}_{\tn{O-Cat}}$, $\cat{W}_{\tn{O-MnCat}}$, $\cat{W}_{\tn{O-TrCat}}$, and $\cat{W}_{\tn{O-Opd}}$ each with a 2-ary morphism that can be drawn as above.%
\footnote{
There is a potentially-confusing level-shift in that there is an operad $\cat{W}_{\tn{O-Opd}}$ for which algebras are operads with object set $O$. In case it's helpful, the object set for $\cat{W}_{\tn{O-Opd}}$ would be $\List(O)\times O$, because every ``box'' has a list of inputs and an output.
}

A $\cat{W}$-algebra $\cat{C}\colon\cat{W}\to\smset$ assigns to each box $b$ the set $\cat{C}(b)$ of ``fillers'' and to each arrangement $\phi$ the function $\cat{C}(\phi)\colon \cat{C}(b_1)\times\cdots\times \cat{C}(b_K)\to \cat{C}(c)$ sending fillers for little boxes to a filler of the big box.%
\footnote{
By a $\cat{W}$-algebra, we mean an operad functor $\cat{W}\to\smset$, where $\smset$ is the operad (multicategory) underlying the monoidal category $(\smset,1,\times)$.
}

\begin{example}
For a set $O$, let $\ob\cat{W}_{\tn{O-Cat}}\coloneqq O\times O$ and
\[\cat{W}_{\tn{O-Cat}}((o_0',o_1),(o_1',o_2),\ldots(o_{K-1}',o_K);(o_0,o_K'))\coloneqq
\begin{cases}
\{1\}&\tn{ if }o_k=o_k'\quad \forall\; 0\leq k\leq K\\
\varnothing&\tn{ otherwise}
\end{cases}
\]
Thus the first diagram shown in \eqref{eqn.wd_examples} represents a map $\phi\colon (A,B),(B,C)\to(A,C)$ in $\cat{W}_{\tn{O-Cat}}$ for some $O$ with $A,B,C:O$. If $\cat{C}\colon\cat{W}_{\tn{O-Cat}}\to\smset$ is an algebra, then the above diagram shows $f_1:\cat{C}(A,B)$, $f_2:\cat{C}(B,C)$ and $\phi(f_1,f_2):\cat{C}(A,C)$. In other words the diagram $\phi$ represents composition in a category. The category of small categories with object set $O$ is equivalent to the category of $\cat{W}_{\tn{O-Cat}}$-algebras.
\end{example}

A variety of examples of these ideas---wiring-diagram operads whose algebras are structured categories---can be found in the literature. \cite{Spivak.Schultz.Rupel:2016a} treats the traced and compact cases, showing in particular that the operad $\cat{W}_{\tn{TrCat}}$ for traced categories is the operad $1\tn{-}\Cat{Cob}$ of oriented 1-dimensional cobordisms. \cite{fong2019hypergraph} treats hypergraph categories, whose associated operad $\cat{W}_{\tn{HypCat}}$ is given by cospans of finite sets. The symmetric monoidal case~\cite{patterson2021wiring} is in some sense the most basic of the three, but its operad $\cat{W}_{\tn{MnCat}}$ is less easy to describe directly.

An advantage of the operadic perspective is that it allows us to more uniformly treat various sorts of category, including ones that may not have an existing name. For example, 
consider the notion of ``traced monoidal categories without identity''; one might balk at this because they wouldn't know whether the ``yanking'' axiom could so easily just be discarded. But as an operad it is fine: it is the suboperad comprising the left class of a certain factorization system on $1\tn{-}\Cat{Cob}$.

Another advantage is that wiring diagram algebras also easily represent $\cat{V}$-enriched categories for monoidal categories $\cat{V}$, namely by using algebras $\cat{W}\to\cat{V}$ valued in $\cat{V}$. For example, a $\vect$-enriched category with object set $O$ is an opeard functor $\cat{W}_{\tn{O-Cat}}\to\vect$.


One can build \emph{dynamic} versions of the structures in \eqref{eqn.wd_examples}, in which the ``fillers'' of a box are ``morphisms that change in time'', i.e.\ coalgebras.  In \cite{shapiro2022dynamic}, a dynamic category is a category enriched in $\org$; similarly for dynamic monoidal categories or dynamic operads. Given the ideas above, these are just operad functors
\[
\cat{W}_{\tn{O-Cat}}\to\org
\qquad
\cat{W}_{\tn{O-MnCat}}\to\org
\qquad
\cat{W}_{\tn{O-Opd}}\to\org
\] 
for any $O:\smset$. Our next goal then is to recall $\org$; in \cref{ch.potentials} we will talk about our choice of $\cat{W}$.

\subsection{The monoidal bicategory $\org$}\label{sec.org}

We recall the monoidal bicategory $\org$ from~\cite{spivak2021learners}; see also~\cite{shapiro2022dynamic}. The objects of $\org$ are polynomials, and for polynomials $p,q$ the hom-category is
\[\org(p,q)\coloneqq\ihom{p,q}\tn{-}\Cat{Coalg}.\]
Its monoidal structure is induced by the Dirichlet product on polynomials. A morphism $p\to q$ in $\org$ is an $\ihom{p,q}$-coalgebra $(S,\,\beta\colon S\to\ihom{p,q}(S))$, specifying for each $s:S$ an action $\act(s):\poly(p,q)$ and an update $\upd_s\colon\ihom{p,q}[\act(s)]\to S$ assigning a successor state to each direction.

\begin{definition}\label{def.pq_coalg}
The category $\ihom{p,q}\coalg$ has as objects pairs $(S,\beta)$ where $S$ is a set and $\beta\colon S\to\ihom{p,q}\tri S$ is a function, and where a morphism from $(S,\beta)$ to $(S',\beta')$ is a function $f\colon S\to S'$ such that $(\ihom{p,q}\tri f)\circ\beta=\beta'\circ f$. We refer to $S$ as the \emph{state set}.
\end{definition}

Unwinding this using \eqref{eqn.dirichlet}, $\beta=(\act^\beta,\upd^\beta)$ decomposes as an \emph{action} function
\[\act^\beta\colon S\to\poly(p,q)=\ihom{p,q}(1)\]
and, for each state $s:S$, an \emph{update} function
\[\upd^\beta_s\colon\sum_{I:p(1)}q\left[\act^\beta_s(I)\right]\to S.\]
For a state $s:S$ and position $I:p(1)$ we often write $\act^\beta_s\colon p\to q$ and $\upd^\beta_s(I)\colon q[\act^\beta_s(I)]\to S$. We may suppress $\beta$ when clear from context.

\begin{proposition}[\cite{spivak2021learners,shapiro2022dynamic}]\label{prop.org}
The hom-categories $\ihom{p,q}\coalg$ assemble into a monoidal bicategory $\org$ with polynomials as objects, with monoidal structure induced by the Dirichlet product on polynomials.
\end{proposition}

This structure arises from the fact that $\coalg\colon(\poly,\yon,\otimes)\to(\smcat,1,\times)$ is lax monoidal:
\[
1\to\yon\coalg
\qqand
p\coalg\times q\coalg\to(p\otimes q)\coalg.
\]
Thus the structure maps $\yon\to[p,p]$ and $[p,q]\otimes[q,r]\to[p,r]$ induce the identity and composition functors
\[
1\to[p,p]\coalg
\qqand
[p,q]\coalg\times[q,r]\coalg\to[p,r]\coalg.
\]

\subsection{The Para construction}\label{sec.para_general}

We use the following form of the Para construction. If a symmetric monoidal category $\cat{A}$ acts on a category $\cat{D}$, the category $\para{\cat{A}}{\cat{D}}$ has the same objects as $\cat{D}$, and a morphism $x\to y$ is a pair $(a,f)$ with $a:\cat{A}$ and
\[
f\colon a\cdot x\to y
\]
in $\cat{D}$. The identity on $x$ has parameter $I$ and the action unitor $I\cdot x\cong x$. The composite of $(a,f)\colon x\to y$ and $(b,g)\colon y\to z$ has parameter $b\otimes a$, or $a\otimes b$ after the symmetry, and uses
\begin{equation}\label{eqn.para_map}
(b\otimes a)\cdot x\cong b\cdot(a\cdot x)\To{b\cdot f}b\cdot y\To{g}z.
\end{equation}

The Para construction extends to a bicategory: a 2-cell $(a,f)\Rightarrow(a',f')$ in $\para{\cat A}{\cat D}(x,y)$ is a morphism $\alpha\colon a\to a'$ in $\cat A$ such that $f'\circ(\alpha\cdot x)=f$. When $\cat A$ is a groupoid, all 2-cells are invertible, and each hom-category $\para{\cat A}{\cat D}(x,y)$ is a groupoid.

\begin{proposition}[Para functoriality from action squares {\cite[\S 3.6, eq.~(3.59)]{capucci2024actegories}}]\label{prop.para_square}
Suppose $\cat A,\cat A'$ are symmetric monoidal categories acting on $\cat D,\cat D'$, that $F\colon\cat A\to\cat A'$ is a strong monoidal functor, $G\colon\cat D\to\cat D'$ is a functor, and $\theta\colon F(a)\cdot G(x)\to G(a\cdot x)$ is a natural transformation compatible with the action unitors and associators:
\begin{equation}\label{eqn.para_square}
\begin{tikzcd}[ampersand replacement=\&]
\cat A\times\cat D\ar[r,"(\cdot)"]\ar[d,"F\times G"']\ar[dr,phantom,"\theta"description]\& \cat D\ar[d,"G"]\\
\cat A'\times\cat D'\ar[r,"(\cdot')"']\& \cat D'.
\end{tikzcd}
\end{equation}
Then there is an induced functor
\[
\para{\cat A}{\cat D}\;\longrightarrow\;\para{\cat A'}{\cat D'},
\]
strong (resp.\ lax) symmetric monoidal whenever $F,G$ are and $\theta$ is monoidal.
\end{proposition}

\begin{proof}
The hypothesis is a 1-cell in the 2-category $\Cat{Act}$ of actegories \cite[\S 3.6, eq.~(3.59)]{capucci2024actegories}; the induced functor records the action of $\Cat{Para}$ on it.\qedhere
\end{proof}

A common source of action-squares is a weakly commuting square of strong monoidal functors, as in \cref{eqn.pnla_polynomial_square}.

\begin{lemma}[Action squares from functor squares]\label{lem.functor_square_action}
Given strong monoidal functors $L,M,N,L'$ between symmetric monoidal categories and a monoidal natural transformation $\theta$ as on the left, one obtains an action square as on the right
\[
\begin{tikzcd}[ampersand replacement=\&]
\cat A\ar[r,"L"]\ar[d,"M"']\ar[dr,phantom,"\theta"description]\& \cat B\ar[d,"N"]\\
\cat A'\ar[r,"L'"']\& \cat B'
\end{tikzcd}
\quad\leadsto\quad
\begin{tikzcd}[ampersand replacement=\&]
\cat A\times\cat B\ar[r,"(\cdot)"]\ar[d,"M\times N"']\& \cat B\ar[d,"N"]\\
\cat A'\times\cat B'\ar[r,"(\cdot')"']\& \cat B'
\end{tikzcd}
\]
by $a\cdot b\coloneqq L(a)\otimes_{\cat B} b$ and similar for $(\cdot')$.
\end{lemma}

\begin{proof}
The natural transformation is given by
\[
L'(M(a))\otimes N(b)\To{\theta_a\otimes\id}N(L(a))\otimes N(b)\To{\cong}N(L(a)\otimes b).
\]
The first factor is the input $\theta$ tensored with identity; the second is the strong monoidal comparison of $N$ (an iso). Compatibility with action unitors and associators reduces to that of $\theta$ and the strong-monoidal data of $L,L',M,N$.\qedhere
\end{proof}

\subsection{Coalgebras as Para over the store action}\label{sec.org_as_para}

The Para construction over the store action of $\smset_\cong$ on $\poly$ (\cref{def.store_action}) repackages the truncation $\org_{\cong}$ (with only coalgebra isomorphisms as 2-cells) exactly:

\begin{proposition}\label{prop.org_as_para}
There is an isomorphism of categories
\[
\para{\Cat{Set}_{\cong}}{\poly}\To{\cong}\org_\cong
\]
\end{proposition}

\begin{proof}
For each category---$\para{\Cat{Set}_{\cong}}{\poly}$ and $\org_\cong$---the object set is $\ob\poly$ and the mapping groupoids $p\to q$ agree
\begin{align*}
	\para{\Cat{Set}_{\cong}}{\poly}(p,q)&=
  \sum_{S:\smset_\cong}\poly(S\yon^S\otimes p, q)\\&\cong
  \sum_{S:\smset_\cong}\poly(S\yon^S,[p,q])\\&\cong
  \sum_{S:\smset_\cong}\poly(S,[p,q]\tri S)=\org_\cong(p,q)
\end{align*}
by \cref{eqn.dirichlet_hom,prop.coalg_as_poly_map}. It remains to show that the identities and composites agree. 

In both cases the identity is indexed by the monoidal unit $1$, and $1\To{\id}\poly(p,p)=[p,p](1)$ and $\yon\To{\id}[p,p]$ are identified as above. Given composable maps $p\To{(S,\varphi)}q\To{(T,\psi)}r$, in both cases the composite is of the form $(TS,\psi\circ\phi)$, and it is easy to see too are identified as above.
\end{proof}

\begin{corollary}[Cotangent learners]\label{cor.cotangent_learners}
There is a strong symmetric monoidal functor
\begin{equation}\label{eqn.cotangent_learners}
\Psi\colon\para{\pvect}{\mfd}\to\org.
\end{equation}
On objects, $M\mapsto\cot M$. On a morphism $(V,f\colon V\times M\to N)$, $\Psi(V,f)$ is the $\ihom{\cot M,\cot N}$-coalgebra with state set $V$ whose action at $v\in V$ is $\cot f_v\colon\cot M\to\cot N$ (with $f_v(m)\coloneqq f(v,m)$) and whose update on a direction $(m,\nu)$ with $\nu\in T^*_{f(v,m)}N$ is
\begin{equation}\label{eqn.cotangent_learner_update}
\upd_v(m,\nu)=v+\sharp_V\bigl((\partial_Vf)_{(v,m)}^\top\nu\bigr).
\end{equation}
\end{corollary}

\begin{proof}
\Cref{prop.para_square,lem.functor_square_action} applied to \cref{prop.pnla_polynomial} give a strong symmetric monoidal functor $\para{\pvect}{\mfd}\to\para{\smset_\cong}{\poly}$ sending $M\mapsto\cot M$ and $f\colon V\times M\to N$ to the morphism with parameter $|V|$ and map
\[
V\yon^V\otimes\cot M\;\To{\theta_V\otimes\id}\;\cot V\otimes\cot M\;\To{\cong}\;\cot(V\times M)\;\To{\cot f}\;\cot N,
\]
where the second arrow is the strong monoidal comparison of $\cot$ (\cref{prop.cot_monoidal}). Composing with the equivalence $\para{\smset_\cong}{\poly}\cong\org_\cong$ of \cref{prop.org_as_para} gives $\Psi$.

The update formula unwinds $\cot f$'s backward part: $(Df)_{(v,m)}^\top\nu=((\partial_Vf)_{(v,m)}^\top\nu,\,(\partial_Mf)_{(v,m)}^\top\nu)\in T_v^*V\oplus T_m^*M$; the $V$-component $\alpha\coloneqq(\partial_Vf)_{(v,m)}^\top\nu$ is identified with the new state $v+\sharp_V(\alpha)$ via $\cot V\cong V\yon^V$ (\cref{prop.pnla_polynomial}).
\end{proof}

\subsection{Example: deep learning as a dynamic monoidal category}\label{sec.deep_learning}

We recast the deep-learning example of~\cite{shapiro2022dynamic} as the restriction of $\Psi$ \eqref{eqn.cotangent_learners} to Euclidean spaces. Take the object set $O\coloneqq\nn$ and assign to each $n\in\nn$ the polynomial
\[
\Psi(\rr^n)=\cot{\rr^n}\cong{\cot{\rr}}\tpow n.
\]
A layer $m\to n$ is a morphism $\rr^m\to\rr^n$ in $\para{\pvect}{\mfd}$: a parameter space $V=\rr^k$, regarded as a paired vector space with $\sharp_V=-\eta\cdot\id$ for a learning rate $\eta>0$, together with a smooth architecture
\[
f\colon V\times\rr^m\to\rr^n.
\]

Applying $\Psi$ gives an $\ihom{\cot{\rr^m},\cot{\rr^n}}$-coalgebra with state set $V$. At $v\in V$, its action is
\[
\cot{f_v}\colon\cot{\rr^m}\to\cot{\rr^n},
\]
whose forward part is $x\mapsto f_v(x)$ and whose backward part is $(Df_v)_x^\top$.

The update at $v$ on a direction $(x,\nu)\in\rr^m\times T^*_{f_v(x)}\rr^n$ is
\[
\upd_v(x,\nu)=v+\sharp_V\bigl((\partial_Vf)_{(v,x)}^\top\nu\bigr)=v-\eta\cdot(\partial_Vf)_{(v,x)}^\top\nu,
\]
by \eqref{eqn.cotangent_learner_update}. When $\nu=dL|_{f_v(x)}$ for a loss $L\colon\rr^n\to\rr$, the chain rule gives $(\partial_Vf)_{(v,x)}^\top\nu=d_v(L\circ f(-,x))$; under the Euclidean identification, $\upd_v$ is one step of stochastic gradient descent.

Since $\Psi$ is a functor, composition in $\org$ reproduces backpropagation. Since $\Psi$ is strong monoidal, parallel composition of layers is transported by the comparison $\cot{\rr^m}\otimes\cot{\rr^{m'}}\cong\cot{\rr^{m+m'}}$. Thus the Euclidean restriction of $\Psi$ assembles into an operad functor
\[\cat{L}\colon\cat{W}_{\nn\tn{-MnCat}}\to\org.\]

\chapter{Potentialized lenses}\label{ch.potentials}

In this chapter we define a symmetric monoidal category $\potlens$ of potentialized lenses. \Cref{sec.lenses_general} develops the abstract construction in one pass: lenses, the backward-monad construction, the represented functor, evaluation at a $\Fun{T}$-algebra, and parameterization. \Cref{sec.potlens_manifolds} specializes the source-side constructions to manifolds: manifold lenses, the potential monad associated to any Lie monoid, and the PN Lie Algebra parameterization, synthesizing them as $\potlens$.

\section{The abstract construction}\label{sec.lenses_general}\label{sec.abstract_representation}

\subsection{Lenses}\label{subsec.lenses}
We first recall the category $\Lens{\cat{C}}$ of lenses in a symmetric monoidal category $(\cat{C},I,\otimes)$; see \cite{hedges2017coherence,spivak2019generalized}. An object is a pair $\lensob{c}$ in which $\outp c$ carries a commutative comonoid structure $(\delta_{\outp c},\varepsilon_{\outp c})$. A morphism $f\colon\lensob{c}\to\lensob{d}$ consists of a pair of maps $f=(\outp f,\inpt f)$ where
\begin{equation}\label{eqn.lens_def}
\outp f\colon \outp c\to\outp d
\qqand
\inpt f\colon \inpt d\otimes\outp c\to\inpt c,
\end{equation}
and $\outp f$ is a comonoid homomorphism.%
\footnote{
When $(\cat{C},1,\times)$ is cartesian monoidal, every object has a unique comonoid structure and every morphism is automatically a comonoid homomorphism~\cite{fox1976coalgebras}.
}
These can be depicted as follows:
\[
\begin{tikzpicture}[oriented WD, bb port length=5pt, bb port sep=1, bb min width=.4cm, bbx=.5cm, bby=.7ex, baseline=(fo)]
  \node[bb={1}{1}] (fo) {$\outp f$};
  \node[bb={1}{2}, above left=6 and 1.5 of fo] (fi) {$\inpt f$};
  \node[circle, draw, fill=black, inner sep=0, minimum size=3pt, left=1 of fo] (dot) {};
  \coordinate (co) at ($(dot)+(-3,0)$);
  \coordinate (do) at ($(fo.east)+(1,0)$);
  \coordinate (ci) at (co |- fi_in1);
  \coordinate (di) at (do |- fi_out1);
  \node[left=0 of co] {$\outp c$};
  \node[right=0 of do] {$\outp d$};
  \node[left=0 of ci] {$\inpt c$};
  \node[right=0 of di] {$\inpt d$};
  \draw[ar] (co) to (dot);
  \draw[ar] (dot) to (fo_in1);
  \draw[ar] (fo_out1) to (do);
  \draw[ar] (fi_in1) to (ci);
  \draw[ar] (di) to (fi_out1);
  \draw[ar] (dot) to[out=30, in=0] (fi_out2);
\end{tikzpicture}
\]
Morphisms compose as follows: given also $g\colon \lensob{d}\to\lensob{e}$, the composite $f\then g$ is given by
\[
\outp{(f\then g)}=\outp f\then\outp g
\qqand
\inpt{(f\then g)}=
	(\inpt{e}\otimes\delta_{\outp{c}})\then
	(\inpt{e}\otimes\outp{f}\otimes\outp{c})\then
	(\inpt{g}\otimes\outp{c})\then
	\inpt{f}.
\]
This dense notation can be depicted in a string-diagram as follows:
\begin{equation}\label{eqn.lens_pic}
\begin{tikzpicture}[oriented WD, bb port length=5pt, bb port sep=1, bb min width=.4cm, bbx=.5cm, bby=.7ex, baseline=(fo)]
  \node[bb={1}{1}] (fo) {$\outp f$};
  \node[bb={1}{1}, right=5 of fo] (go) {$\outp g$\vphantom{$\inpt f$}};
  \node[bb={1}{2}, above left=6 and 1.5 of fo] (fi) {$\inpt f$};
  \node[bb={1}{2}, above left=6 and 1.5 of go] (gi) {$\inpt g$\vphantom{$\inpt f$}};
  \node[circle, draw, fill=black, inner sep=0, minimum size=3pt, left=1 of fo] (dot1) {};
  \node[circle, draw, fill=black, inner sep=0, minimum size=3pt] at ($(fo.east)!.75!(go.west)$) (dot2) {};
  \coordinate (co) at ($(dot1)+(-3,0)$);
  \coordinate (eo) at ($(go.east)+(1,0)$);
  \coordinate (ci) at (co |- fi_in1);
  \coordinate (ei) at (eo |- gi_out1);
  \node[left=0 of co] {$\outp c$};
  \node[right=0 of eo] {$\outp e$};
  \node[left=0 of ci] {$\inpt c$};
  \node[right=0 of ei] {$\inpt e$};
  \draw[ar] (co) to (dot1);
  \draw[ar] (dot1) to (fo_in1);
  \draw[ar] (fo_out1) to node[above, font=\scriptsize] {$\outp d$} (dot2);
  \draw[ar] (dot2) to (go_in1);
  \draw[ar] (go_out1) to (eo);
  \draw[ar] (fi_in1) to (ci);
  \draw[ar] (ei) to (gi_out1);
  \draw[ar] (gi_in1) to[out=180, in=0] node[above, font=\scriptsize] {$\inpt d$} (fi_out1);
  \draw[ar] (dot1) to[out=30, in=0] (fi_out2);
  \draw[ar] (dot2) to[out=30, in=0] (gi_out2);
\end{tikzpicture}
\end{equation}
The identity on $\lensob{c}$ is $(\id_{\outp c},\id_{\inpt c}\otimes\varepsilon_{\outp c})$, i.e.\ counit in the second argument. The category $\Lens{\cat{C}}$ inherits a monoidal structure with unit $\binom{I}{I}$ and
\[
\lensob{c}\otimes\lensob{d}\coloneqq\lensobs{c}{d}.
\]



\begin{proposition}\label{prop.Theta}
For any symmetric monoidal closed category $(\cat{C},I,\otimes)$, there is a lax monoidal functor
\[\Theta\colon\Lens{\cat{C}}\to\cat{C},\qquad
\lensob{c}\mapsto\outp c\otimes\ihom{\inpt c,I}.\]
\end{proposition}
\begin{proof}
Given a lens $f\colon\lensob c\to\lensob d$, define $\Theta(f)$ as the composite
\[
\begin{tikzcd}[column sep=45pt, row sep=30pt]
  \outp c\otimes[\inpt c,I]\ar[r, "{\delta_{\outp c}\otimes\id}"]\ar[d, dashed, "\Theta(f)"']&
  \outp c\otimes\outp c\otimes[\inpt c,I]\ar[r, "{\outp f\otimes\coev}"]&
  \outp d\otimes[\inpt d,\inpt d\otimes\outp c\otimes[\inpt c,I]]\ar[d, "{\id\otimes[\inpt d,\inpt f\otimes\id]}"]\\
  \outp d\otimes[\inpt d,I]&&\outp d\otimes[\inpt d,\inpt c\otimes[\inpt c,I]]\ar[ll, "{\id\otimes[\inpt d,\ev]}"]
\end{tikzcd}
\]
One checks that this preserves identity and composition by the fact that $\outp{f}$ is a comonoid homomorphism and the closure triangle identity $[\inpt d,\ev]\circ\coev=\id_{[\inpt d,I]}$.

To show it is (normal) lax, we have $I\otimes I\cong I$ and $[\inpt c,I]\otimes[\inpt d,I]\to[\inpt c\otimes\inpt d,I]$.\qedhere
\end{proof}



\subsection{Backward monads on lenses}\label{subsec.backwards}

We will add a notion of potentials to our lenses by proving a more general result: any strong monad on $\cat{C}$ induces a comonad on $\Lens{\cat{C}}$ by applying it to the input / backward-pass part. We pun on the word ``backwards'' to refer to the fact that it applies to the backward pass and the fact that it turns a monad into a comonad.

Recall that a monad $\Fun{T}$ on a monoidal category is \emph{strong} if there are natural maps $\sigma\colon X\otimes\Fun{T}Y\to\Fun{T}(X\otimes Y)$, satisfying four diagrams~\cite{kock1972strong}. If moreover $X$ carries a counit $\varepsilon_X\colon X\to I$, then naturality of $\sigma$ in its first argument together with the unit-coherence axiom of strength gives a \emph{counital projection} $\sigma_{X,Y}\then\Fun{T}(\varepsilon_X\otimes\id_Y)=\varepsilon_X\otimes\id_{\Fun{T}Y}$:
\[
\begin{tikzcd}
X\otimes\Fun{T}Y\ar[r, "\sigma"]\ar[d, "\varepsilon_X\otimes\id"']&
\Fun{T}(X\otimes Y)\ar[d, "\Fun{T}(\varepsilon_X\otimes\id)"]\\
\Fun{T}Y\ar[r, equal]&
\Fun{T}Y
\end{tikzcd}
\]
When $\cat{C}$ is cartesian, this specializes to the familiar projection formula $\sigma\then\Fun{T}\pi_Y=\pi_{\Fun{T}Y}$.


\begin{proposition}
Let $(\Fun{T},\eta,\mu,\sigma)$ be a strong monad on a symmetric monoidal category $(\cat{C},I,\otimes)$. Then there is a comonad $\Lens{\Fun{T}}$ on $\Lens{\cat{C}}$ sending $\lensob{c}\mapsto \binom{\Fun{T}\inpt c}{\outp c}$ and sending a map $(\outp f,\inpt f)$ to $(\outp f, \sigma\then\Fun{T}\inpt f)$
\[
\outp c\otimes\Fun{T}\inpt d\To{\sigma}\Fun{T}(\outp c\otimes\inpt d)\To{\Fun{T}\inpt f}\Fun{T}\inpt c
\]
\end{proposition}
\begin{proof}
The counit and comultiplication of $\Lens{\Fun{T}}$ are given by the unit and multiplication of the monad, together with the counit $\varepsilon_{\outp c}$:
\[
\outp c\otimes\inpt c\To{\varepsilon_{\outp c}\otimes\id}\inpt c\To{\eta}\Fun{T}\inpt c
\qqand
\outp c\otimes\Fun{T}\Fun{T}\inpt c\To{\varepsilon_{\outp c}\otimes\id}\Fun{T}\Fun{T}\inpt c\To{\mu}\Fun{T}\inpt c
\]
We need to prove that $\Lens{\Fun{T}}$ preserves identity and composition (is a functor) and that the functor is counital and coassociative. These four axioms correspond precisely to the four axioms of a monad strength.
%The unit axiom~\cite[(1.5)]{kock1972strong} (left unitor coherence) gives preservation of identity, via $\sigma_{X,Y}\then\Fun{T}\pi_Y=\pi_{\Fun{T}Y}$. The $\eta$-coherence axiom of the strong monad gives the two counit laws. The $\mu$-coherence axiom, together with associativity of $\mu$, gives coassociativity.

We include only the least trivial here: preservation of composition. Writing $f_o,g_o$ for forward components and $f_i\colon\outp c\otimes\inpt d\to\inpt c$, $g_i\colon\outp d\otimes\inpt e\to\inpt d$ for backward components, we must show that the following commutes (after appending $\Fun{T}f_i\colon\Fun{T}(\outp c\otimes\inpt d)\to\Fun{T}\inpt c$ on the right):
\[
\begin{tikzcd}[column sep=1.6em, row sep=3em, every label/.append style={font=\scriptsize}, font=\small, ampersand replacement=\&]
\outp c\otimes\Fun{T}\inpt e
  \ar[r, "\delta\otimes\id"]
  \ar[d, "\sigma"']
\& (\outp c)^{\otimes 2}\otimes\Fun{T}\inpt e
  \ar[r, "\id\otimes f_o\otimes\id"]
  \ar[d, "\sigma"]
\&[5pt] \outp c\otimes\outp d\otimes\Fun{T}\inpt e
  \ar[r, "\id\otimes\sigma"]
  \ar[dr, "\sigma" description]
\& \outp c\otimes\Fun{T}(\outp d\otimes\inpt e)
  \ar[r, "\id\otimes\Fun{T}g_i"]
  \ar[d, "\sigma"]
\& \outp c\otimes\Fun{T}\inpt d
  \ar[d, "\sigma"]
\\
\Fun{T}(\outp c\otimes\inpt e)
  \ar[r, "\Fun{T}(\delta\otimes\id)"']
\& \Fun{T}((\outp c)^{\otimes 2}\otimes\inpt e)
  \ar[rr, "\Fun{T}(\id\otimes f_o\otimes\id)"']
\&
\& \Fun{T}(\outp c\otimes\outp d\otimes\inpt e)
  \ar[r, "\Fun{T}(\id\otimes g_i)"']
\& \Fun{T}(\outp c\otimes\inpt d)
\end{tikzcd}
\]
The first two squares are naturality of $\sigma$ in its first argument; the middle triangle is the associator axiom of strength~\cite[(1.6)]{kock1972strong}, which merges the nested $\sigma$'s; the last square is naturality of $\sigma$ in its second argument.
\end{proof}

We denote the coKleisli category of $\Lens{\Fun{T}}$ by $\Lens{\cat{C}}^{\Fun{T}}$. A map in $\Lens{\cat{C}}^{\Fun{T}}$ is a map of the form $\binom{\Fun{T}\inpt c}{\outp c}\to\lensob{d}$ in $\Lens{\cat{C}}$, i.e.
\[
\outp f\colon \outp c\to\outp d
\qqand
\inpt f\colon \inpt d\otimes\outp c\to\Fun{T}\inpt c.
\]

Recall that monoidal (equivalently, commutative) monads $(\Fun{T},\tau)$ have a natural strength
\[
\sigma_{X,Y}\colon X\otimes\Fun{T}Y\To{\eta_X\otimes\id}\Fun{T}X\otimes\Fun{T}Y\To{\tau_{X,Y}}\Fun{T}(X\otimes Y).
\]
and that monoidal monads $\Fun{T}$ preserve monoids. The following generalizes to any $\otimes$-monoid in place of $I$ (\cref{rmk.general_z}).

\begin{proposition}\label{prop.Theta_T_alpha}
Suppose $(\Fun{T},\tau)$ is a monoidal monad and that the unit carries an algebra structure $\alpha\colon\Fun{T}I\to I$,
\[
\begin{tikzcd}[ampersand replacement=\&]
I\ar[r,"\eta_I"]\ar[dr,equal]\& \Fun{T}I\ar[d,"\alpha"]\\
\& I
\end{tikzcd}
\qquad
\begin{tikzcd}[ampersand replacement=\&]
\Fun{T}\Fun{T}I\ar[r,"\mu_I"]\ar[d,"\Fun{T}\alpha"']\& \Fun{T}I\ar[d,"\alpha"]\\
\Fun{T}I\ar[r,"\alpha"']\& I
\end{tikzcd}
\]
such that $\alpha$ is also a monoid homomorphism:
\[
\begin{tikzcd}[ampersand replacement=\&]
I\ar[r,"\tau_0"]\ar[dr, equal]\& \Fun{T}I\ar[d,"\alpha"]\\
\& I
\end{tikzcd}
\qquad
\begin{tikzcd}[ampersand replacement=\&]
\Fun{T}I\otimes\Fun{T}I\ar[r,"\tau_{I,I}"]\ar[d,"\alpha\otimes\alpha"']\& \Fun{T}(I\otimes I)\ar[r, equal]\& \Fun{T}I\ar[d,"\alpha"]\\
I\otimes I\ar[rr, equal]\&\& I
\end{tikzcd}
\]
Then $\Theta$ from \cref{prop.Theta} extends to a lax monoidal functor
\[
\Theta^{\Fun{T},\alpha}\colon\Lens{\cat{C}}^\Fun{T}\to\cat{C},\qquad
\lensob{c}\mapsto\outp c\otimes\ihom{\inpt c,I}.
\qedhere
\]
\end{proposition}
\begin{proof}
Given a coKleisli morphism $f\colon\lensob{c}\coto\lensob{d}$, define $\Theta^{\Fun{T},\alpha}(f)$ as the composite
\[
\begin{tikzcd}[column sep=45pt, row sep=22pt, ampersand replacement=\&]
\outp c\otimes[\inpt c,I]\ar[r,"{\delta_{\outp c}\otimes\id}"]\ar[dd, dashed, "\Theta^{\Fun{T},\alpha}(f)"']
\& \outp c\otimes\outp c\otimes[\inpt c,I]\ar[r,"\outp f\otimes\coev"]
\& \outp d\otimes[\inpt d,\inpt d\otimes\outp c\otimes[\inpt c,I]]\ar[d,"{\id\otimes[\inpt d,\inpt f\otimes\id]}"]
\\
\&
\& \outp d\otimes[\inpt d,\Fun{T}\inpt c\otimes[\inpt c,I]]\ar[d,"{\id\otimes[\inpt d,\sigma]}"]
\\
\outp d\otimes[\inpt d,I]
\& \outp d\otimes[\inpt d,\Fun{T}I]\ar[l,"{\id\otimes[\inpt d,\alpha]}"]
\& \outp d\otimes[\inpt d,\Fun{T}(\inpt c\otimes[\inpt c,I])]\ar[l,"{\id\otimes[\inpt d,\Fun{T}\ev]}"]
\end{tikzcd}
\]
Everything else is the same as in \cref{prop.Theta}.
\end{proof}


\begin{remark}\label{rmk.general_z}
The construction above uses the monoidal unit $z\coloneqq I$ as the target of evaluation. More generally, for any $\otimes$-monoid $(z,e_z,m_z)$ equipped with a $\Fun{T}$-algebra $\alpha\colon\Fun{T}z\to z$ which is also a monoid homomorphism, the above generalizes to a lax monoidal functor  
\[
\Theta^{\Fun{T},z,\alpha}\colon\Lens{\cat{C}}^\Fun{T}\to\cat{C},\qquad
\lensob{c}\mapsto\outp c\otimes\ihom{\inpt c,z}.
\qedhere
\]
\end{remark}

\subsection{Parameterized lenses}\label{sec.abstract_parameters}

We apply the Para construction (\cref{sec.para_general}) to potentialized lenses. A strong monoidal functor $J\colon\cat{A}\to\cat{C}$ landing in commutative comonoids gives an $\cat{A}$-action on $\Lens{\cat{C}}^\Fun{T}$ by $a\cdot\lensob c\coloneqq\binom{\inpt c}{J(a)\otimes\outp c}$ (the extra $J(a)$ in the outp slot is killed in the backward leg by its counit), and $\Theta^{\Fun{T},\alpha}$ intertwines this with the action $a\cdot x\coloneqq J(a)\otimes x$ on $\cat{C}$. Applying \cref{prop.para_square,lem.functor_square_action} gives a strong symmetric monoidal functor
\[
\para{\cat{A}}{\Lens{\cat{C}}^\Fun{T}}\to\para{\cat{A}}{\cat{C}},
\]
sending a parameterized coKleisli lens $(a,f\colon a\cdot\lensob c\coto\lensob d)$ to the parameterized $\cat C$-morphism $J(a)\otimes\Theta^{\Fun{T},\alpha}\lensob c\to\Theta^{\Fun{T},\alpha}\lensob d$.

\section{Potentialized manifold lenses}\label{sec.potlens_manifolds}

We specialize the abstract construction of \cref{sec.lenses_general} to $\cat C=\mfd$ (cartesian monoidal). The resulting category $\lmfd\coloneqq\Lens\mfd$ of \emph{manifold lenses} has pairs of manifolds $\lensob M$ as objects and pairs $f=(\outp f,\inpt f)$ of smooth maps as morphisms (\eqref{eqn.lens_def}).

Any Lie monoid $(G,\cdot,e)$ gives a strong monad on $\mfd$:
\[
\Fun T_G(X)\coloneqq X\times G,\qquad\eta_X(x)\coloneqq(x,e),\qquad\mu_X(x,g,g')\coloneqq(x,g\cdot g'),
\]
with strength $X\times(Y\times G)\To{\cong}(X\times Y)\times G$ by reassociation. The backward-monad construction (\cref{subsec.backwards}) produces a comonad $\Lens{\Fun T_G}$ on $\lmfd$; we write $\lmfd^G$ for its coKleisli category. A morphism $\lensob L\coto\lensob M$ in $\lmfd^G$ has backward component
\[
\inpt f\colon\outp L\times\inpt M\to\inpt L\times G,
\]
a usual backward map augmented by a smooth $G$-valued map; coKleisli composition multiplies $G$-components.

The case of primary interest is $G=(\rr,0,+)$, the \emph{potential monad}: the $\rr$-component
\[
U\colon\outp L\times\inpt M\to\rr
\]
is the \emph{potential}, and coKleisli composition adds potentials.

We lift the $\pvect$-action on $\mfd$ (\cref{prop.pnla_smc_action}) to $\lmfd$ by acting on the $\outp M$-slot,
\[
(\cdot)\colon\pvect\times\lmfd\to\lmfd,\qquad V\cdot\lensob M\coloneqq\binom{\inpt M}{V\times\outp M},
\]
and define the category of \emph{potentialized lenses}
\[
\potlens\coloneqq\para{\pvect}{\lmfd^\rr}.
\]
A morphism $f\colon\lensob L\to\lensob M$ in $\potlens$ is a tuple $(V,\outp f,\inpt f,U)$ with $V$ a paired vector space,
\begin{equation}\label{eqn.map_paralens}
\outp f\colon V\times\outp L\to\outp M
\qqand
\inpt f\colon V\times\outp L\times\inpt M\to\inpt L,
\end{equation}
and a smooth potential $U\colon V\times\outp L\times\inpt M\to\rr$:
\[
\begin{tikzpicture}[oriented WD, bb port length=5pt, bb port sep=1, bb min width=.4cm, bbx=.5cm, bby=.7ex, baseline=(fo)]
  \node[bb={2}{1}] (fo) {$\outp f$};
  \node[bb={1}{3}, above left=4 and 4 of fo] (fi) {$\inpt f$};
  \node[bb={1}{3}, above right=3 and 0 of fi] (V) {$U$};
  \node[circle, draw, fill=black, inner sep=0, minimum size=3pt, left=1 of fo_in2] (dotL) {};
  \coordinate (Lo) at ($(dotL)+(-4.5,0)$);
  \coordinate (Mo) at ($(fo.east)+(1,0)$);
  \coordinate (Li) at (Lo |- fi_in1);
  \coordinate (Mi) at (Mo |- fi_out3);
  \node[circle, draw, fill=black, inner sep=0, minimum size=3pt] at ($(fi_out3)+(2,0)$) (dotMI) {};
  \node[left=0 of Lo] {$\outp L$};
  \node[right=0 of Mo] {$\outp M$};
  \node[left=0 of Li] {$\inpt L$};
  \node[right=0 of Mi] {$\inpt M$};
  \draw[ar] (Lo) to (dotL);
  \draw[ar] (dotL) to (fo_in2);
  \draw[ar] (fo_out1) to (Mo);
  \draw[ar] (fi_in1) to (Li);
  \draw[ar] (Mi) to (dotMI);
  \draw[ar] (dotMI) to (fi_out3);
  \draw[ar] (dotL) to[out=60, in=0] (fi_out2);
  \draw[ar] (dotMI) to[out=60, in=0] (V_out3);
  \draw[ar] (dotL) to[out=60, in=0] (V_out2);
  %
  \node [above left=1 and 0 of Lo] (P) {$V$};
  \node[circle, draw, fill=black, inner sep=0, minimum size=3pt] at (dotL|-P) (dotP) {};
  \draw[ar] (P) to (dotP);
  \draw[ar] (dotP) to (fo_in1);
  \draw[ar] (dotP) to[out=60, in=0] (fi_out1);
  \draw[ar] (dotP) to[out=60, in=0] (V_out1);
  \coordinate (R) at ($(V.west)+(-2,0)$);
  \node[left=0 of R] {$\rr$};
  \draw[ar] (V_in1) to (R);
\end{tikzpicture}
\]
The identity on $\lensob{L}$ has parameter $\rr^0:\pvect$, backward projection on $\lmfd$, and $U=0$ (the unit $e=0$ of $\rr$). Given also $g=(W,\outp g,\inpt g,U')\colon\lensob{M}\to\lensob{N}$, the composite $f\then g$ has parameter $V\times W$, inherits forward and backward components from $\lmfd$, and its potential adds the two via the coKleisli multiplication $\mu=+$:
\begin{align*}
U\bigl(v,\ell,\,\inpt g(w,\outp f(v,\ell),n')\bigr)+U'\bigl(w,\outp f(v,\ell),n'\bigr).
\end{align*}
The monoidal product is inherited from $\lmfd$. The following depicts composition, though note that we include it only for the algorithmically-minded reader's convenience: \emph{it follows from the abstract framework we already have developed}:
\[
\begin{tikzpicture}[oriented WD, bb port length=5pt, bb port sep=1, bb min width=.4cm, bbx=.5cm, bby=.7ex, baseline=(fo)]
  \node[bb={2}{1}] (fo) {$\outp f$};
  \node[bb={2}{1}, right=7 of fo] (go) {$\outp g$};
  \node[bb={1}{3}, above left=4 and 3 of fo] (fi) {$\inpt f$};
  \node[bb={1}{3}, above left=4 and 3 of go] (gi) {$\inpt g$};
  \node[bb={1}{3}, above right=1 and 0 of fi] (V) {$U$};
  \node[bb={1}{3}, above right=1 and 0 of gi] (W) {$U'$};
  \node[bb={1}{2}, above left=-1 and 2 of V] (plus) {$+$};
  %
  \node[circle, draw, fill=black, inner sep=0, minimum size=3pt, left=1 of fo_in2] (dotL) {};
  \node[circle, draw, fill=black, inner sep=0, minimum size=3pt] at ($(gi.east|-fo.east)+(1,0)$) (dotM) {};
  %
  \coordinate (Lo) at ($(dotL)+(-5,0)$);
  \coordinate (No) at ($(go.east)+(1,0)$);
  \coordinate (Li) at (Lo |- fi_in1);
  \coordinate (Ni) at (No |- gi_out3);
  \node[left=0 of Lo] {$\outp L$};
  \node[right=0 of No] {$\outp N$};
  \node[left=0 of Li] {$\inpt L$};
  \node[right=0 of Ni] {$\inpt N$};
  %
  \draw[ar] (Lo) to (dotL);
  \draw[ar] (dotL) to (fo_in2);
  \draw[ar] (fo_out1) to node[above, font=\scriptsize] {$\outp M$} (dotM);
  \draw[ar] (dotM) to (go_in2);
  \draw[ar] (go_out1) to (No);
  \draw[ar] (fi_in1) to (Li);
  \node[circle, draw, fill=black, inner sep=0, minimum size=3pt] at ($(gi_out3)+(3,0)$) (dotNI) {};
  \draw[ar] (Ni) to (dotNI);
  \draw[ar] (dotNI) to (gi_out3);
  \node[circle, draw, fill=black, inner sep=0, minimum size=3pt] at ($(fi_out3)+(3,0)$) (dotMI) {};
  \draw[ar] (gi_in1) to[out=180, in=0] node[above, font=\scriptsize] {$\inpt M$} (dotMI);
  \draw[ar] (dotMI) to (fi_out3);
  \draw[ar] (dotL) to[out=60, in=0] (fi_out2);
  \draw[ar] (dotM) to[out=60, in=0] (gi_out2);
  %
  \node[above left=1 and 0 of Lo] (P) {$V$};
  \node[above left=4 and 0 of Lo] (Q) {$W$};
  \node[circle, draw, fill=black, inner sep=0, minimum size=3pt] at (dotL|-P) (dotP) {};
  \node[circle, draw, fill=black, inner sep=0, minimum size=3pt] at (dotM|-Q) (dotQ) {};
  \draw[ar] (P) to (dotP);
  \draw[ar] (Q) to (dotQ);
  \draw[ar] (dotP) to (fo_in1);
  \draw[ar] (dotP) to[out=60, in=0] (fi_out1);
  \draw[ar] (dotQ) to (go_in1);
  \draw[ar] (dotQ) to[out=60, in=0] (gi_out1);
  %
  \draw[ar] (dotP) to[out=60, in=0] (V_out1);
  \draw[ar] (dotL) to[out=60, in=0] (V_out2);
  \draw[ar] (dotMI) to[out=120, in=0] (V_out3);
  \draw[ar] (dotQ) to[out=60, in=0] (W_out1);
  \draw[ar] (dotM) to[out=60, in=0] (W_out2);
  \draw[ar] (dotNI) to[out=120, in=0] (W_out3);
  %
  \draw[ar] (V_in1) to[out=180, in=0] (plus_out2);
  \draw[ar] (W_in1) to[out=180, in=0] (plus_out1);
  \coordinate (R) at ($(plus.west)+(-1,0)$);
  \node[left=0 of R] {$\rr$};
  \draw[ar] (plus_in1) to (R);
\end{tikzpicture}
\]


\section{Lenses represent serial, parallel, feedback, and splitting}\label{sec.plenspot_ops}

Recall from \cref{sec.dynamic_structures} that operads---and hence monoidal categories---can be used to encode compositional syntax. In fact, the monoidal category $\Lens{\cat{M}}$ of lenses in a symmetric monoidal category $\cat{M}$ with a supply of commutative comonoids does just that for ``traced symmetric monoidal semi-categories equipped with a supply of comonoids'', as we now explain.

By a \emph{semi-category} we mean almost a category, except no identities (like a semigroup is a group without identities). So our goal is to show that $\Lens{C}$ interprets diagrams that look like this:
\begin{equation}\label{eqn.arb_wd}
\varphi\coloneqq
\begin{tikzpicture}[oriented WD, bb min width=.6cm, bb port sep=1, bbx=.6cm, bby=1ex, bb port length=2.5pt, baseline=(X2)]
  \node[bb port sep=2, bb={2}{2}] (X1) {};
  \node[bb={1}{1}, right=.7 of X1_out1] (X2) {};
  \node[bb={0}{0}, fit={(X1) (X2) ($(X1.north)+(0,2.5)$)}] (Y) {};
  \node[coordinate] at (Y.west|-X1_in2) (Y_in1) {};
  \node[coordinate] at (Y.east|-X2_out1) (Y_out1) {};
  \node[coordinate] at (Y.east|-X1_out2) (Y_out2) {};
  \draw[ar] (Y_in1) -- node[above, font=\scriptsize] {$A$} (X1_in2);
  \draw[ar] (X1_out1) -- node[above, font=\scriptsize] {$B$} (X2_in1);
  \draw[ar] (X2_out1) -- node[above, font=\scriptsize] {$C$} (Y_out1);
  \draw[ar] (X1_out2) -- node[below, font=\scriptsize] {$D$} (Y_out2);
  \draw[ar] let \p1=(X2.north east), \p2=(X1.north west), \n1={\y2+\bby}, \n2=\bbportlen in
          (X2_out1) to[in=0,in looseness=.9] (\x1+.7*\n2,\n1) -- node[above, font=\scriptsize] {$C$} (\x2-.7*\n2,\n1) to[out=180] (X1_in1);
  \end{tikzpicture}
  \end{equation}

Looking at the two inner boxes and the outer box, we see that this diagram should be represented in $\Lens{\cat{M}}$ by a map of the form
\[
\varphi\colon\binom{C\times A}{B\times D}\otimes\binom{B}{C}\to\binom{A}{C\times D}
\]
Unfolding, we need two maps in $\cat{M}$:
\[
B\times D\times C\to C\times D
\qqand
B\times D\times C\times A\to C\times A\times B
\]
But these maps are obvious: they are just projections.

Lenses are much more general than wiring diagrams. Indeed, all of the maps required by wiring diagrams are projections and diagonals, whereas there are many more maps, e.g.\ of the form $B\times D\times C\to C\times D$, than the projection(s). Looking at \eqref{eqn.arb_wd}, one can imagine that the outputs of the boxes, $B,D,C$ can be combined in any way in order to get the output of the big box (similarly for the inputs of the little boxes), not just by the direct channel shown in the pictures.

But adding parameterizations as in \cref{sec.para_general} makes things even more general: the parameter can affect how information is sent around:
\[
V\times B\times D\times C\to C\times D
\qqand
V\times B\times D\times C\times A\to C\times A\times B.
\]
In other words, the wiring itself can be smoothly varied. And finally by adding in a potential $U$ as in \cref{sec.potlens_manifolds}, we can grade the information flow.


\chapter{From lenses to dynamics}\label{ch.potlens_to_org}

This chapter defines the functor $\Phi\colon\potlens\to\org$ by applying the abstract construction of \cref{sec.abstract_representation} to polynomial functors. The key infrastructure --- that $\Cat{Para}$ over the store action of $\Cat{Set}_{\cong}$ on $\poly$ presents $\org$ --- was set up in \cref{sec.org_as_para,cor.cotangent_learners}. The remaining inputs are the cotangent functor, the potential monad, and the PN Lie Algebra parameterization.

\section{Specialization to polynomial functors}\label{sec.specialize_source}

Take $\cat C=\poly$, with the strong monad
\[
\Fun T(p)\coloneqq p\otimes\cot\rr,
\]
whose unit and multiplication are $\cot 0\colon\yon\to\cot\rr$ and $\cot{+}\colon\cot\rr\otimes\cot\rr\to\cot\rr$ from \cref{prop.cot_hopf}. The algebra $\alpha\colon\cot\rr\to\yon$ required by \cref{prop.Theta_T_alpha} must be a monoid homomorphism from $\cot\rr$ (a Hopf monoid by \cref{prop.cot_hopf}) to the trivial monoid on $\yon$.

\begin{lemma}\label{lem.alpha_constant}
The $\Fun T$-algebra structures on $\yon$ that preserve the trivial monoid are in bijection with $\rr$: each such $\alpha\colon\cot\rr\to\yon$ has direction-action at each $x\in\rr$ given by a constant $c\in\rr$ independent of $x$, and every constant $c$ arises this way.
\end{lemma}

\begin{proof}
The position-action of $\alpha$ is the unique map $\rr\to 1$, so $\alpha$ is determined by its direction-action, a family $c_x\coloneqq\bk\alpha_x(*)\colon T_x^*\rr\cong\rr$ indexed by $x\in\rr$. Evaluating the monoid-preservation law $\cot{+}\then\alpha=(\alpha\otimes\alpha)\then\lambda_\yon$ on the unique direction of $\yon$ at position $(x,y)\in\rr\times\rr$ gives
\[
(c_{x+y},\,c_{x+y})=(c_x,\,c_y)\colon T_x^*\rr\oplus T_y^*\rr,
\]
using that $D(+)=(1,1)$ transposes to the diagonal. Setting $x=0$ yields $c_y=c_0$ for all $y$, so $c_y$ is a constant $c\coloneqq c_0\in\rr$. Conversely, for any constant $c$, setting $\bk\alpha_x(*)\coloneqq c$ defines $\alpha$, and the algebra and monoid-hom laws all reduce to the trivial $(c,c)$-computation; the unit laws hold because $T^*_*\rr^0$ is a one-element set.
\end{proof}

We choose
\[
\dec\colon\cot\rr\to\yon,\qquad\bk\dec_x(*)\coloneqq -1.
\]
The sign fixes the sign of Hamilton's equations below; $c=0$ gives the trivial algebra (no potential), $c=+1$ the opposite-sign physics. \Cref{rmk.general_z} allows algebras valued in monoids $z\neq\yon$; we do not pursue them.

The lensification of $\cot\colon\mfd\to\poly$ (\cref{prop.cot_monoidal,lem.cot_comonoid}) and $\Theta^{\Fun T,\dec}\colon\Lens\poly^{\Fun T}\to\poly$ (\cref{prop.Theta_T_alpha}) compose to
\[
\Phi_1\colon\lmfd^\rr\to\poly,\qquad
\Phi_1\lensob M=\cot{\outp M}\otimes\ihom{\cot{\inpt M},\yon}.
\]
Specializing the parameterized-lens construction (\cref{sec.abstract_parameters}) to $\cat A=\pvect$ with $J=\cot$ (strong monoidal and comonoid-valued by \cref{prop.cot_monoidal,prop.pnla_smc_action,lem.cot_comonoid}) gives
\[
\potlens=\para{\pvect}{\lmfd^\rr}\to\para{\pvect}{\poly},
\]
where $\pvect$ acts on $\poly$ by $V\cdot p\coloneqq\cot V\otimes p$. A morphism $f\colon\lensob L\to\lensob M$ in $\potlens$ with parameter $(V,\sharp_V)$ gives the map
\begin{equation}\label{eqn.parameter_to_action}
\cot V\otimes\Phi_1\lensob L\to\Phi_1\lensob M.
\end{equation}

\section{The cotangent-state refinement}\label{sec.cotangent_state_refinement}

For the dynamics in this paper, we want $T^*V$ rather than $V$ as the state space in $\org$. This is achieved by changing the action of $\pvect$ on $\poly$ from $V\cdot p=\cot V\otimes p$ to $V\cdot p=\cot{T^*V}\otimes p$, mediated by a (non-iso) natural transformation $\rho_V\colon\cot{T^*V}\to\cot V$.

For any paired vector space $(V,\sharp_V)$, $T^*V\cong V\oplus V^*$ by \eqref{eqn.tangent_vec}. We write $(u,\pi)$ for a \emph{position} $u\in V$ and \emph{momentum} $\pi\in V^*$. The cotangent space decomposes as
\[
T^*_{(u,\pi)}(T^*V)\cong T^*_uV\oplus T^*_\pi(V^*)\cong V^*\oplus V,
\]
again by \eqref{eqn.tangent_vec}, using $(V^*)^*\cong V$ on the $V^*$-factor.

\begin{definition}[Pairing projection]\label{def.pairing_proj}
The \emph{pairing projection} $\rho_V\colon\cot{T^*V}\to\cot V$ has position-action $(u,\pi)\mapsto u$ and direction-action at $(u,\pi)$
\[
T^*_uV\cong V^*\To{\sigma\mapsto(\sigma,\,\sharp_V(\pi))}V^*\oplus V\cong T^*_{(u,\pi)}(T^*V).
\]
\end{definition}

On positions, $\rho_V$ forgets momentum; on directions, it passes the incoming covector $\sigma$ through to the $V^*$-slot and injects the pairing-lifted momentum $\sharp_V(\pi)\in V$ into the $V$-slot, independently of $\sigma$. So $\rho_V$ is not of the form $\cot h$ for any smooth $h\colon T^*V\to V$, and not invertible. It is monoidal and natural in $\pvect$-isomorphisms because $\sharp_{V_1\oplus V_2}=\sharp_{V_1}\oplus\sharp_{V_2}$ (\cref{prop.TT_monoidal}).

\Cref{lem.functor_square_action} applied with $L=\cot\circ\Fun{inc}$, $L'=\cot\circ\Fun{inc}\circ T^*$, $M=N=\id$, and $\theta_V=\rho_V$ produces an action square comparing the actions $V\cdot p=\cot V\otimes p$ and $V\cdot p=\cot{T^*V}\otimes p$ on $\poly$. \Cref{prop.para_square} gives a functor between the two Para constructions, and \cref{cor.cotangent_learners}'s coalgebra interpretation, applied at $T^*V$ (using \cref{prop.pnla_polynomial}: $\cot{T^*V}\cong (T^*V)\yon^{T^*V}$ since $T^*V$ is a paired vector space by \cref{prop.canonical_symplectic_pairing}), packages the result as an $\org$-morphism with state set $T^*V$. This is the 1-cell $\Phi(f)$ in $\org$.

\begin{remark}[Tangent-state alternative]\label{rmk.tangent_state}
Taking $T$ instead of $T^*$ (\cref{prop.TT_endofunctors}) fits the same framework. The analogous \emph{tangent projection} $\tilde\rho_V\colon\cot{TV}\to\cot V$ has position-action $(u,X)\mapsto u$ and direction-action $\sigma\mapsto(\sigma,\,\flat_V(X))$, using the flat map $\flat_V\colon V\to V^*$ to place $X$ in the second summand of $T^*_{(u,X)}(TV)\cong V^*\oplus V^*$. This produces a functor $\potlens\to\org$ with state set $TV$ (Lagrangian-style, second-order), which we don't develop here; the preference for $T^*$ is physics (cotangent bundles carry the canonical symplectic pairing of \cref{prop.canonical_symplectic_pairing}, recovering Hamilton's equations).
\end{remark}

\begin{remark}[Legendre equivalence of Hamiltonian and Lagrangian $\Phi$]\label{rmk.legendre_equivalence}
By \cref{prop.legendre}, $\sharp\colon T^*\Rightarrow T$ is a monoidal natural isomorphism, so $\rho$ and $\tilde\rho$ correspond via $\cot\sharp$. The two functors $\Phi^{T^*},\Phi^{T}\colon\potlens\to\org$ are naturally isomorphic: $\sharp_V\colon T^*V\to TV$ induces a 2-cell in $\org$ identifying $\Phi^{T^*}(f)$ with $\Phi^{T}(f)$. This is the \emph{Legendre equivalence} of Hamiltonian and Lagrangian mechanics in our setting.
\end{remark}

\section{Functoriality and lax monoidality}\label{sec.functoriality}

\begin{theorem}\label{thm.functor}
$\Phi\colon\potlens\to\org$ is a lax symmetric monoidal functor.
\end{theorem}

\begin{proof}
$\Phi$ is the composite
\[
\potlens=\para{\pvect}{\lmfd^\rr}\To{\para\pvect{\Phi_1}}\para{\pvect}{\poly}\To{\rho^*}\para{\pvect,\,\cot\circ T^*}{\poly}\To{(*)}\org,
\]
where $\Phi_1$ is from \cref{sec.specialize_source}, $\rho^*$ is the action-square refinement from \cref{sec.cotangent_state_refinement}, and $(*)$ is the coalgebra interpretation at state $T^*V$ (\cref{prop.coalg_as_poly_map,prop.pnla_polynomial}). Each is functorial; their composite is $\Phi$.

For lax monoidality: $\Phi_1$ is lax (composing strong $\widehat\cot$ with lax $\Theta^{\Fun T,\dec}$ from \cref{prop.Theta_T_alpha}), $\para\pvect{\Phi_1}$ inherits laxness, $\rho^*$ is monoidal (\cref{prop.para_square}, since $\rho$ is monoidal), and $(*)$ is strong by \cref{prop.pnla_polynomial,prop.cot_monoidal}. Composing yields lax. The morphism laxator is the closedness pairing $\ihom{\cot{\inpt L},\yon}\otimes\ihom{\cot{\inpt{L'}},\yon}\to\ihom{\cot{\inpt L\times\inpt{L'}},\yon}$ of \cref{prop.Theta}, which is not invertible (it misses covector fields on $\inpt L\times\inpt{L'}$ with genuine mixed dependence), so $\Phi$ is lax, not strong.\qedhere
\end{proof}

\begin{proposition}\label{prop.potlens_monoidal}
$\potlens$ inherits from $\lmfd$ and $\pvect$ a symmetric monoidal structure with unit $\binom{\rr^0}{\rr^0}$ and tensor $\lensob L\otimes\lensob{L'}\coloneqq\lensobs{L}{L'}$. On morphisms, $f\otimes f'$ has parameter $V\oplus V'$ (with pairing $\sharp_V\oplus\sharp_{V'}$), independent forward and backward components, and potential $U^\otimes=U\circ\pr+U'\circ\pr'$.
\end{proposition}

\begin{proof}
Para over a symmetric monoidal action of an SMC on an SMC is itself symmetric monoidal. Here $\pvect$ is SMC (\cref{prop.pnla_monoidal}) and $\lmfd^\rr$ is SMC via lensification of $\mfd$ together with the writer monad on $\rr$; the $\pvect$-action on $\lmfd^\rr$ is symmetric monoidal via the forgetful $\pvect\to\mfd$ (\cref{prop.pnla_smc_action}). The morphism formula reads off from the Para tensor: parameter $V\oplus V'$, slotwise forward/backward, additive potential.\qedhere
\end{proof}

\begin{remark}[Lagrangian formulation]\label{rmk.lagrangian}
The Lagrangian formulation is the same dynamics written on $TV$ rather than on $T^*V$. A point of $T^*V$ is a position together with a momentum, while a point of $TV$ is a position together with a velocity. The component $\sharp_V\colon T^*V\To{\cong}TV$ of \cref{prop.legendre} converts momentum to velocity, with inverse $\flat_V$, so $X_H$ transports to a vector field on $TV$:
\[
Y_H\colon
TV\To{\flat_V}T^*V\To{X_H}T(T^*V)\To{T(\sharp_V)}T(TV).
\]
In coordinates $(q,w)\in TV$ with $\alpha\coloneqq\flat_V(w)$, the integral curves of $Y_H$ satisfy
\begin{equation}\label{eqn.lagrangian_pushforward}
\dot q=\partial_\alpha H|_{(q,\alpha)},\qquad \dot w=-\sharp_V\partial_v H|_{(q,\alpha)}.
\end{equation}
For an integral curve $\gamma(t)=(q(t),w(t))$ in $TV$, its tangent is a point
\[
\dot\gamma(t)=((q,w),(\dot q,\dot w))\in T(TV)\cong TV\oplus TV.
\]
Thus a vector field on $TV$ is a second-order equation for $q$ exactly when its first tangent component is $w$.

When the pairing on $V$ is symmetric and the Hamiltonian $H\colon T^*V\to\rr$ is of the form
\[
H(v,\alpha)=\tfrac{1}{2}\alpha(\sharp_V\alpha)+U(v),
\]
where $U\colon V\to\rr$ (the \emph{potential}), then
\[
d_\alpha\left(\tfrac{1}{2}\alpha(\sharp_V\alpha)\right)(\beta)
=\tfrac{1}{2}\bigl(\beta(\sharp_V\alpha)+\alpha(\sharp_V\beta)\bigr)
=\beta(\sharp_V\alpha)
\]
for every $\beta\in V^*$. Thus, under $(V^*)^*\cong V$, $\partial_\alpha H=\sharp_V\alpha$, while $\partial_v H=dU$, so \eqref{eqn.lagrangian_pushforward} becomes
\[
\dot q=w,\qquad \dot w=-\sharp_V dU(q).
\]
The first equation identifies $w$ with the velocity of $q$; the second says $\ddot q=\dot w=-\sharp_V dU(q)$, equivalently $\flat_V(\ddot q)=-dU(q)$. This is Newton's second law in this notation: $\flat_V$ converts acceleration to the corresponding force covector, and $-dU(q)$ is the force determined by the potential.

The associated Lagrangian function is obtained from the same Hamiltonian by rewriting momentum as velocity:
\[
L(q,w)\coloneqq \flat_V(w)(w)-H(q,\flat_V(w))
=\tfrac{1}{2}\flat_V(w)(w)-U(q).
\]
The Euler--Lagrange equation for this $L$,
\[
\frac{d}{dt}\partial_w L=\partial_q L,
\]
is exactly $\flat_V(\ddot q)=-dU(q)$, so the Lagrangian has recovered the same second-order equation without using momentum as a state variable.
In this special case one often says ``kinetic plus potential'' for $H$ and ``kinetic minus potential'' for $L$; these names only describe this separable example. The main point is the change of state space: Hamiltonian dynamics uses $(q,\alpha)\in T^*V$, while Lagrangian dynamics uses $(q,w)\in TV$, and $\sharp_V$ identifies the two descriptions.\qedhere
\end{remark}

%--------------------------------------------------------------------
\chapter{Example: the spring cell}\label{ch.spring}
%--------------------------------------------------------------------

We illustrate $\Phi$ with a one-dimensional spring cell and recover the discrete and continuum wave equations by wiring copies of it.

\begin{definition}\label{def.spring_cell}
Fix $\kappa,\mu>0$. The \emph{spring cell} is the $\potlens$-morphism
\[
f_\mathrm{spring}\colon\binom{\rr^2}{\rr}\to\binom{\rr^2}{\rr}
\]
with input $\rr^2=\{(\ell,r)\}$ (left and right neighbor positions), output $\rr$ (cell's own position), paired vector space parameter $V\coloneqq(\rr,0,+)$ with pairing $\sharp_V(p)\coloneqq p/\mu$, and data
\[
\outp f(x,y)\coloneqq x,
\quad
\inpt f(x,y,(\ell,r))\coloneqq(\ell,r),
\quad
U(x,y,(\ell,r))\coloneqq\tfrac{\kappa}{2}(x-\ell)^2+\tfrac{\kappa}{2}(r-x)^2.
\]
\end{definition}

\begin{remark}\label{rmk.spring_pairing}
The pairing $\sharp_V(p)=p/\mu$ is chosen, not the standard Euclidean pairing $\sharp_V(p)=p$. The free parameter $\mu$ encodes inverse mass; the standard pairing corresponds to $\mu=1$.
\end{remark}

\begin{proposition}\label{prop.spring_dynamics}
At state $(x,p): T^*V$ on input $(y,\alpha)$ with vanishing received covector $\nu=0$ and input direction $(\ell,r)\in \rr^2$, the update of $\Phi(f_\mathrm{spring})$ is
\[
\dot x = p/\mu,
\qquad
\dot p = \kappa(\ell+r-2x).
\]
\end{proposition}

\begin{proof}
Unwinding $\Phi(f_\mathrm{spring})$ at $(x,p)$: the $\sharp_V(p)=p/\mu$ component of $\dot x$ is the $V$-slot of the pairing projection (\cref{def.pairing_proj}). The $\dot p$ component receives contributions from $U$, $\outp f$, and $\inpt f$ along the chain induced by $\rho_V$; only the $U$-contribution is nonzero here ($\partial_V\inpt f=0$ and $\nu=0$), giving $\dot p=-\partial_xU=\kappa(\ell+r-2x)$.\qedhere
\end{proof}

\begin{remark}[Harmonic oscillator Hamiltonian]\label{rmk.spring_hamiltonian}
The spring cell's state-update is textbook Hamilton's equations for the harmonic-oscillator Hamiltonian
\[
H(x,p) \;=\; \frac{p^2}{2\mu} \;+\; \tfrac{\kappa}{2}(x-\ell)^2 \;+\; \tfrac{\kappa}{2}(r-x)^2
\]
on $T^*\rr\cong\rr^2$: kinetic energy $p^2/(2\mu)$ (mass $\mu$) plus two Hooke potentials (one spring to the left neighbor at $\ell$, one to the right neighbor at $r$, both with spring constant $\kappa$). \Cref{prop.spring_dynamics}'s equations $\dot x=p/\mu$, $\dot p=\kappa(\ell+r-2x)$ are precisely $\dot x=\partial_p H$, $\dot p=-\partial_x H$. Differentiating once more recovers Newton's second law $\mu\ddot x=\kappa(\ell+r-2x)$.
\end{remark}

\begin{proposition}\label{prop.spring_wave}
Wiring a $\zz$-indexed chain of spring cells with each cell's input $(\ell,r)$ identified with the outputs of its left and right neighbors yields joint dynamics
\[
\mu\ddot x_i = \kappa(x_{i-1}-2x_i+x_{i+1}).
\]
In the continuum limit $x_i\to u(i\Delta x,t)$ with density $\rho\coloneqq\mu/\Delta x$ and tension $T\coloneqq\kappa\Delta x$, the chain approaches the wave equation $\rho\, u_{tt}=T\,u_{xx}$.
\end{proposition}

\begin{proof}
The wiring is a closed-system 1-cell in $\org$ obtained by $\otimes$-product and feedback identifying each cell's input slot with neighbors' output slots (cf.\ \cref{sec.plenspot_ops}). Apply \cref{prop.spring_dynamics} cell-wise and time-differentiate $\dot x_i=p_i/\mu$; substitute $\dot p_i=\kappa(x_{i-1}+x_{i+1}-2x_i)$ to obtain the discrete equation, and pass to the continuum limit.
\end{proof}


\printbibliography

\end{document}
